\chapter{前沿展望——LLM 的下一步}
\label{ch:outlook}
%% ============================================================

\section{推理模型的进化：从技巧到系统能力}

推理（reasoning）能力的发展经历了几个清晰的阶段：

\textbf{第一阶段：提示工程}。
Chain-of-thought（CoT）提示的发现表明，
只需在 prompt 中加入 “Let's think step by step”，
模型的数学和逻辑推理能力就会显著提升。
这说明模型内部已经具备了推理能力，
只是需要被正确地「激活」。

\textbf{第二阶段：推理微调}。
通过在包含推理链的数据上进行 SFT，
让模型默认生成推理过程。
这比提示工程更稳定，模型不再依赖用户提供正确的 prompt。

\textbf{第三阶段：专用推理模型}。
OpenAI o1/o3/o4-mini、DeepSeek-R1~\cite{deepseekr1}、
Claude 的 extended thinking 模式
代表了推理能力的系统化训练
（其中 o1/o3 系列和 DeepSeek-R1 是专门训练的推理模型，
而 Claude 的 extended thinking 是通用模型的推理增强模式，
二者技术路线有本质区别）。
这些模型通过 RL（如 GRPO~\cite{shao2024grpo}）
学会了在推理过程中进行自我验证、回溯和错误修正。
关键创新是将推理链的生成视为一个可以用 RL 优化的过程——
模型不仅学会了「思考」，还学会了「思考如何思考」。

\subsection{推理模型的技术路线分化}

截至 2026 年初，推理模型已经形成了三条明显不同的技术路线：

\textbf{路线一：专用推理模型（OpenAI o 系列）}。
OpenAI 的 o1（2024 年 9 月）→ o3（2024 年 12 月预览，2025 年 4 月正式发布）
→ o4-mini（2025 年 4 月）
形成了一条清晰的演进路径。
这些模型在训练时就被专门优化为「深度思考者」，
它们生成长而详尽的推理链，
在数学（AIME 2025：o3 达到 96.7\%）、
编程（Codeforces rating $>$2700）和科学推理（GPQA Diamond：87.7\%）
等需要深度推理的任务上表现突出。
o4-mini 是这条路线的效率优化版本：
以显著更低的成本和延迟提供接近 o3 的推理能力，
在许多评估中甚至略微超越 o3。
这表明推理能力的 scaling 不仅可以通过增大模型实现，
还可以通过更好的训练方法在更小的模型上实现。

\textbf{路线二：开源推理模型（DeepSeek-R1）}。
DeepSeek-R1~\cite{deepseekr1} 用一个关键实验
揭示了推理能力的涌现机制：
在 R1-Zero 实验中，
他们仅用格式化的 RL reward（不使用任何 SFT 数据），
就让 base model 自发学会了长链推理和自我反思。
这个发现意义深远——
它表明推理能力不需要人工标注的推理链数据来训练，
RL 本身就能诱导出这种能力。
R1 的完整版本在 AIME 2024 上达到 79.8\%
（与 o1 正式版的 83.3\% 接近），
并以 MIT 许可证开源，
催生了大量基于 R1 蒸馏的小型推理模型
（如 R1-Distill-Qwen-32B、R1-Distill-LLaMA-70B 等）。

\textbf{路线三：混合推理模式（Claude extended thinking）}。
Anthropic 的 Claude 3.7 Sonnet（2025 年 2 月）
开创了「混合推理」的概念，
同一个模型可以在快速回答和深度推理之间切换。
用户无需选择不同的模型，
只需启用 extended thinking 模式，
模型就会在回答前生成可见的思考过程。
Claude Opus 4.6 进一步引入了可配置的 thinking budget
（最高 128K tokens 的推理预算），
用户可以根据任务复杂度和成本考量做出选择。
这种路线的优势是灵活性——
简单问题无需支付推理 token 的成本，
复杂问题可以投入更多推理计算。

\textbf{其他值得关注的推理模型}：
Google 的 Gemini 2.5 Pro（内置 “thinking” 模式，
2025 年 3 月在 LMSYS Chatbot Arena 登顶）；
Grok 3（xAI，2025 年 2 月，内置 “Big Brain” 推理模式）；
Qwen-QwQ（阿里巴巴，开源的推理专用模型）。

\subsection{2026 年春季：中国推理模型的「集体爆发」}

2025 年 12 月至 2026 年 2 月，中国 AI 实验室掀起了一波前所未有的推理模型发布潮。
路透社将其描述为“抢先竞争”（preemptive competition），
六家中国头部实验室在三周内密集发布旗舰模型，
形成了 2026 年春季的“模型发布海啸”。

\textbf{DeepSeek V3.2/V3.2-Speciale}（2025 年 12 月）
将竞赛级推理能力推向了新高度：
国际数学奥林匹克（IMO）获得金牌（35/42 分），
国际大学生程序设计竞赛（ICPC）取得第二名，
国际信息学奥林匹克（IOI）第十名，
AIME 2025 达到 96.0\%，超越了 GPT-5-High 的表现。
这些成绩标志着 LLM 的推理能力已经在最严苛的竞赛环境中
达到了人类顶尖选手的水平。

\textbf{DeepSeek V4}（预计 2026 年 4 月发布）
是 DeepSeek 的下一代旗舰模型，目前仍在开发中。
据接近 DeepSeek 的消息源透露，V4 预计将实现多项突破：
万亿参数级 MoE 架构（约 1T 总参数，$\sim$320B 活跃参数），
原生多模态（文本/图像/视频/音频的统一理解与生成），
100 万 token 上下文窗口，
以及基于 Engram 论文的长期记忆（LTM）机制。
特别值得关注的是，V4 据报道针对华为昇腾 910B/C 进行了专门优化，
并在预发布阶段刻意未向 NVIDIA 和 AMD 提供预览访问，
这被视为中国 AI 芯片独立路线的重要信号。

\textbf{豆包 Seed 2.0 Pro}（2026 年 2 月 14 日，字节跳动）
在数学推理和代码能力上全面超越同期竞品：
AIME 2025 达到 98.3\%（创当时最高纪录），
Codeforces rating 达到 3020（超越 99.9\% 的人类选手），
SWE-bench Verified 76.5\%。
Seed 2.0 Pro 的发布证明了字节跳动
在大规模基础模型训练上的快速追赶能力。

\textbf{Kimi K2.5}（2026 年 1 月 27 日，月之暗面）
是 Kimi 系列首个万亿参数模型，
采用 1T 总参数、384 个专家（8 个活跃/token）的 MoE 架构。
其最大创新是 \textbf{Agent Swarm} 系统，
通过 PARL（Parallel Agent Reinforcement Learning）协议
协调多达 100 个子 Agent 并行执行任务，
最长可运行 1500 步，实现 80\% 的运行时间缩减。
AIME 2025 达到 96.1\%，HMMT 达到 95.4\%。
Kimi K2.5 同时也是 Kimi 系列首个多模态旗舰，
搭载了 MoonViT-3D（400M 视觉编码器）。

\textbf{阶跃星辰 Step 3.5 Flash}（2026 年 2 月 2 日）
在极致效率上做出了突破：
196B 总参数、11B 活跃参数的 MoE 架构，
可以在单块 RTX 4090（24GB 显存）上运行。
AIME 2025 达到 97.3\%：
以不到 GPT-5 百分之一的推理成本
实现了接近 GPT-5 的推理性能。
Step AI 的 CTO 曾坦言，
DeepSeek R1 的推理范式对其构成了“降维打击”，
促使他们从纯 scaling 路线转向效率优化。

\begin{corebox}[2026 年春季模型发布潮的竞争动态]
2025 年 12 月至 2026 年 2 月的密集发布潮
不仅仅是技术进步的自然结果，
更是竞争策略的产物。
路透社的报道证实了这种“抢先竞争”的逻辑：
各实验室争相在 DeepSeek V4 正式发布前
占领话语权和用户心智。
六家头部实验室（DeepSeek、字节跳动、月之暗面、阶跃星辰、智谱、阿里巴巴）
在三周内全部发布了旗舰模型，
几乎所有模型都以开放权重发布，
并伴随着激进的定价策略。
这种竞争强度在全球 AI 领域前所未有。
\end{corebox}

\begin{whybox}[推理时计算 vs.\ 训练时计算的经济学]
一个深刻的趋势转变正在发生：
传统的 Scaling Laws 关注训练时计算——
投入更多算力训练更大的模型。
但推理时计算（test-time compute）提供了另一条路径，
保持模型大小不变，通过在推理时投入更多计算
（更长的思考链、更多的候选方案探索）来提升输出质量。
这改变了成本结构：训练成本是一次性的，
而推理成本随使用量线性增长。
对于使用量大的应用，
一个更小但「想得更久」的模型
可能比一个更大但「直接回答」的模型更具经济优势。
o4-mini 的成功正是这一趋势的体现，
它以 o3 约 1/5 的成本提供了相当甚至更好的推理能力。
\end{whybox}

\section{多模态：走向统一理解}

2025--2026 年，
纯文本 LLM 正在快速转向多模态模型，
同时理解文本、图像、音频甚至视频。
Gemini 3、GPT-4o、Claude 4.5 Sonnet 都是原生多模态的。
多模态不再是「附加功能」，而是「基础能力」。

\subsection{视觉-语言模型的架构演进}

视觉-语言模型（Vision-Language Model, VLM）
的发展经历了几代架构：

\textbf{CLIP 范式}：
对比学习，训练一个图像编码器和一个文本编码器，
使匹配的图文对在表示空间中靠近。
CLIP 提供了强大的视觉表示，但不能生成文本。

\textbf{LLaVA 范式}：
视觉编码器 + 投影层 + LLM。
将预训练的视觉编码器（如 ViT/SigLIP）的输出
通过一个简单的投影层（线性层或 MLP）
映射到 LLM 的输入空间，
然后与文本 token 一起送入 LLM。
这种架构简洁有效，
是当前大多数开源 VLM 的基础。

\textbf{原生多模态}：
GPT-4o 和 Gemini 3（基于 Gemini 架构~\cite{anil2023gemini} 的后续版本）代表了下一步，
不再将视觉作为「插件」，
而是在预训练阶段就联合训练视觉和语言理解。
这种方法成本更高（需要海量的图文配对数据），
但产生的模型对视觉信息的理解更深入、更自然。

\subsection{音频：从转录到原生理解}

音频能力的发展也遵循类似的路径：

\textbf{转录范式}：使用 Whisper 等专用模型将语音转为文本，
再用 LLM 处理文本。这是目前最稳定的方案，
但丢失了语音中的非文本信息（语调、情感、停顿）。

\textbf{原生音频}：GPT-4o 的实时语音模式
和 Gemini 的原生音频输入
代表了音频理解的下一步，
模型直接处理音频波形，
不经过中间的文本转录步骤。
这使得模型能够理解语调变化、背景声音、
甚至多人对话中的说话者切换。

实时语音交互的技术挑战巨大：
模型需要在数百毫秒内开始生成回复（否则对话会感觉卡顿），
同时能够被用户打断（“barge-in” 检测），
还需要处理环境噪声和远场拾音。

\textbf{音频生成}同样在快速演进。
从早期的 TTS（文本转语音）管线
，先生成文本再用 Vocoder 合成语音，
到 GPT-4o 的原生音频生成（直接从内部表示产生语音 token），
模型生成的语音在自然度、情感表达和多语言支持上
都有了质的飞跃。
2025--2026 年的一个重要趋势是\textbf{语音克隆与风格迁移}——
模型可以基于几秒钟的参考音频
生成具有特定说话风格的语音输出，
这在客服、有声书和个性化助手场景中有巨大价值，
但同时引发了关于 deepfake 音频和身份冒充的严重安全担忧。

\textbf{音乐生成}是音频 AI 的另一个前沿。
\textbf{Suno}（2023 年底发布，持续迭代至 2025 年底 v4）
和 \textbf{Udio}（2024 年发布）
展示了 AI 可以生成包含歌词、旋律、和声和编曲的
完整音乐作品。
Suno v4 生成的歌曲在盲测中
被部分听众评价为“与专业制作的音乐难以区分”，
这对音乐产业的影响可能堪比摄影技术对绘画的冲击。
然而，版权问题仍然是音乐生成 AI 面临的最大法律挑战，
多家唱片公司已对 Suno 和 Udio 提起版权侵权诉讼，
核心争议在于：
如果 AI 在受版权保护的音乐上训练，
它生成的作品是否构成对原作的“衍生”。

\textbf{声音理解}也在快速演进，
超越了传统的语音识别范畴。
Google 的 Gemini 3 和 OpenAI 的 GPT-4o
不仅能转录语音内容，
还能识别环境声音
（“这段录音中有鸟叫声和远处的交通噪音”）、
分析音乐特征
（“这首曲子是 A 小调，4/4 拍，BPM 约 120”）、
甚至理解非语言的声音信号
（“这个人的声音听起来焦虑”）。
这种\textbf{通用声音理解}能力的出现
使得音频模态不再只是“语音转文字的管道”，
而是与视觉和文本并列的
独立信息来源。

\subsection{视频理解与世界模型}

视频是多模态中计算成本最高的模态。
一段 1 分钟的 30fps 视频包含 1800 帧：
如果每帧都编码为数百个 token，
模型需要处理数十万个视觉 token。

当前的处理策略包括：
\textbf{关键帧采样}：从视频中均匀或自适应地选取少量帧；
\textbf{视频编码器}——使用 3D 卷积或时序 Transformer
将连续帧压缩为更紧凑的表示；
\textbf{分层处理}——先用轻量级模型定位关键片段，
再用大模型深入分析。

长视频理解（数十分钟到数小时）
仍然是一个开放问题，
它同时挑战了模型的上下文窗口、
视觉理解能力和长期记忆能力。

\textbf{世界模型（World Models）}
代表了视频理解的更深层方向，
不是简单地生成“看起来像视频的像素”，
而是试图学习世界运行的底层规则。

\subsubsection{Sora 系列：从视频生成到世界模拟}

OpenAI 的 Sora 最初于 2024 年 2 月以技术报告形式亮相，
核心架构是 \textbf{Diffusion Transformer（DiT）}~\cite{peebles2023dit}：
将 Transformer 的注意力机制与扩散模型的去噪训练相结合。
具体来说，Sora 将视频分解为时空 patch
（每个 patch 覆盖若干帧的一个空间区域），
通过一个视频压缩网络将原始像素压缩到低维的潜在空间，
然后在潜在空间中使用 DiT 进行自回归或双向去噪。
这种设计的关键优势是\textbf{分辨率和时长的灵活性}——
因为 Transformer 天然处理可变长度的 token 序列，
Sora 可以生成不同分辨率和宽高比的视频，
而不需要为每种配置训练单独的模型。

Sora 的训练数据包括大量的网络视频和人工标注的高质量片段。
OpenAI 在技术报告中声称，
Sora 展现了一些令人惊讶的“涌现”物理理解，
物体的碰撞、反射、流体运动等现象
在生成的视频中呈现出一定的物理合理性，
尽管模型从未被显式地教授物理定律。
这引发了一个深刻的争论：
\textbf{视频生成模型是否在隐式地学习物理？}

\textbf{Sora 2}（2025 年 9 月发布）
在初代基础上做了重大升级：
支持 1080p 视频生成（初代仅为 720p），
最长 1 分钟的连续视频（初代约 10--20 秒），
并加入了\textbf{同步音频生成}能力，
模型可以为生成的视频自动配上环境音效和背景音乐。
Sora 2 的另一个重要进步是\textbf{角色一致性}：
在多个连续镜头中保持同一角色的外观一致，
这是初代 Sora 的主要弱点之一。

\subsubsection{世界模型的理论之争}

“视频生成即世界模拟”的论点
引发了 AI 理论界最激烈的辩论之一。

\textbf{扩散模型路线}认为，
如果视频生成模型在足够多的真实世界视频上训练，
它将隐式地学会物理规律，
因为生成物理上合理的视频
\textbf{是}最有效的数据压缩策略。
这种观点的支持者指出，
Sora 等模型在没有任何物理引擎辅助的情况下
就展现出了一定的物理直觉，
这暗示着“足够大的模型 + 足够多的视频数据
= 隐式的物理引擎”。

\textbf{JEPA（Joint Embedding Predictive Architecture）路线}
是 Yann LeCun 在 Meta FAIR 倡导的替代方案~\cite{lecun2022path}。
LeCun 长期批评自回归和扩散模型的“像素级预测”范式，
认为在像素空间中预测未来帧是极端低效的，
因为未来的不确定性在像素层面是巨大的
（一片树叶可能向左飘也可能向右飘），
但在\textbf{语义层面}是有限的
（树叶会落地，不会飞上天）。
JEPA 的核心思想是在\textbf{学到的抽象表示空间}中进行预测，
而不是在原始像素空间中。
V-JEPA~\cite{bardes2024vjepa}（2024 年）
展示了这种方法在视频理解任务上的潜力，
它通过掩码预测（masking）在视频的潜在空间中学习表示，
不需要像素级的重建损失。
V-JEPA 2~\cite{vjepa2_2025}（2025 年）
进一步将这一范式推向了更大的规模，
在多个视频理解基准上取得了 SOTA。

两种路线的分歧本质上是关于\textbf{世界模型应该建模什么}：
扩散模型建模的是“世界看起来怎样”（像素分布），
JEPA 建模的是“世界的状态是什么”（语义表示）。
截至 2026 年初，扩散模型在视觉质量上占据压倒性优势
（Sora、Kling、Runway 生成的视频视觉上令人惊叹），
但 JEPA 在下游任务（如机器人规划、因果推理）
上可能更有潜力，
因为机器人不需要知道每个像素的颜色，
只需要知道“桌子上有一个杯子，杯子可能会被碰倒”。

\subsubsection{视频生成竞赛}

Sora 引发了全球范围的视频生成竞赛。
截至 2026 年初，主要竞争者包括：

\textbf{Kling}（快手，2024 年 6 月发布，持续迭代至 2025 年底）
是中国视频生成的旗舰。
Kling 2.0 支持 1080p、最长 2 分钟的视频生成，
在运动连贯性和物理合理性上与 Sora 2 不相上下，
但在中文场景理解（文化特定的视觉内容）上更有优势。
Kling 的核心技术包括自研的 3D VAE 视频编码器
和基于 DiT 的时空注意力架构。

\textbf{Runway Gen-3 Alpha}（2024 年 6 月发布）
以创意工作流集成著称，
Runway 将视频生成嵌入专业视频编辑工作流，
支持精细的摄像机运动控制、
风格迁移和多模态条件输入（文本 + 图像 + 草图）。
Gen-3 Alpha Turbo 大幅降低了生成延迟，
使得交互式创作成为可能。

\textbf{Pika 2.0}（2024 年底--2025 年初）
在消费者市场取得了成功，
其“Scene to Video”功能
允许用户通过一张照片生成包含人物动作的短视频，
病毒式传播使 Pika 成为视频生成领域
用户量增长最快的产品之一。

\textbf{Google Veo 2}（2024 年 12 月发布）
在视频生成质量评估中多次超越 Sora，
支持 4K 分辨率和更长时长的生成。
Veo 2 整合了 Google 在图像生成（Imagen 3）
和视频理解（VideoMAE）方面的技术积累，
是 Google“全栈多模态”战略的重要拼图。

\textbf{字节跳动 Seaweed}（2025 年发布）
是字节跳动在视频生成领域的布局，
利用抖音/TikTok 的海量视频数据优势，
在短视频生成场景（15--60 秒、竖屏、社交媒体风格）
上形成了独特的竞争力。

\subsubsection{世界模型的局限性}

然而，当前的世界模型距离真正理解物理定律还有很大差距。
具体的局限性包括：

\textbf{物理不一致性}。
生成的视频仍然频繁出现违反物理定律的场景，
物体凭空出现或消失、
重力方向突然改变、
流体行为不符合流体力学、
镜面反射与光源位置不一致。
这些错误暴露了模型学到的是
“统计上的视觉模式”而非“因果性的物理规律”。

\textbf{长期连贯性}。
随着生成视频时长的增加，
角色外观、场景布局和叙事逻辑的一致性急剧下降。
目前即使是最先进的模型
也很难生成超过 2 分钟的物理上连贯的视频。

\textbf{可控性}。
用户很难精确控制生成内容的细节，
“让角色走到桌子前，拿起红色杯子，转身离开”
这样的精细指令对当前模型来说仍然是巨大挑战。

尽管如此，世界模型的方向对于
机器人、自动驾驶和游戏 AI 至关重要：
一个能够预测物理世界演变的模型
可以用于规划和决策。
自动驾驶场景中，
世界模型可以“想象”一辆前方车辆突然变道后的场景，
帮助规划系统提前做出应对决策~\cite{hu2023gaia1}。

\subsection{Omni-模型：全模态输入与输出}

多模态的终极形态是 \textbf{omni-model}：
不仅能理解多种模态的输入，
还能生成多种模态的输出（文本、图像、音频、视频）。

GPT-4o 的 “o” 就代表 “omni”。
它可以接收文本、图像和音频输入，
并生成文本和音频输出。
GPT-4o 在 2025 年 3 月获得了原生图像生成能力，
不再依赖独立的 DALL-E 模型。
Gemini 3 更进一步，原生支持图像和音频的生成。

统一模型的关键挑战是\textbf{训练数据的模态平衡}——
互联网上的文本数据远多于高质量的图文配对数据，
而高质量的视频-文本配对数据更加稀缺。
如何在预训练中平衡不同模态的数据比例，
避免某种模态的能力被其他模态「挤压」，
是一个活跃的研究方向。

\textbf{Omni-model 的架构路线}正在经历分化。
一种路径是\textbf{统一 tokenizer}：
将所有模态（文本、图像、音频、视频）
映射到同一个离散 token 空间中，
然后用单一的 Transformer 进行自回归生成。
这是最优雅的方案，但对 tokenizer 的设计要求极高，
图像和音频的连续信号如何高效离散化
而不丢失关键信息，仍然是开放问题。
另一种路径是\textbf{模块化生成}——
使用统一的理解模型处理输入，
但为不同模态的输出维护独立的生成头
（text head、image diffusion head、audio vocoder）。
这种方案工程上更易实现，
且允许每个生成模块独立优化，
但增加了模型的参数量和部署复杂度。

\textbf{实时交互}是 omni-model 的杀手级应用场景。
GPT-4o 的实时语音对话
和 Gemini 的 Project Astra 展示了
一个能够同时看、听、说的 AI 助手的潜力：
用户可以对着手机摄像头提问，
AI 在理解视觉和语音输入后，
以自然语音和屏幕标注同时回复。
这种交互模式的延迟要求极其严苛
（端到端延迟需低于 500 毫秒才能维持自然对话节奏），
对推理基础设施提出了远超纯文本模型的挑战。
Qwen3-Omni（2025 年 Q3）进一步将开源模型
推入了 omni-modal 竞争，
它是首个开源的 omni-modal 模型，
支持文本、图像、音频的输入与输出，
在实时交互延迟上与 GPT-4o 形成了直接竞争。

\subsection{2026 年多模态的新突破}

2026 年初，多模态模型出现了几个值得关注的新方向：

\textbf{原生多模态成为标配}。
DeepSeek V4（预计 2026 年 4 月）将是 DeepSeek 首个
原生多模态模型，文本、图像、视频、音频的统一理解与生成
不再通过外挂编码器实现，而是在预训练阶段联合学习。
这标志着中国开源模型在多模态技术路线上
追上了 GPT-4o 和 Gemini 3 的原生多模态范式。

\textbf{Kimi K2.5 的 MoonViT-3D}（2026 年 1 月 27 日）
引入了 400M 参数的 3D 视觉编码器，
这是 Kimi 系列首个多模态旗舰模型。
MoonViT-3D 通过三维时空建模
对视频内容的理解能力显著优于传统的 2D 视觉编码器，
特别是在需要理解运动、因果关系和空间推理的任务上。

\textbf{GLM-Image}（2026 年 1 月 16 日，智谱 AI）
具有里程碑意义，它是首个完全在国产芯片上训练的
SOTA 级多模态模型，
使用华为昇腾 Atlas 800T A2 集群和昇思 MindSpore 框架，
实现了多模态理解和图像生成的双重能力，
且未使用任何 NVIDIA GPU。
这证明了国产芯片生态已经具备支撑前沿多模态训练的能力。

\textbf{OpenRouter Healer Alpha}（2026 年 3 月 11 日）
以匿名形式在 OpenRouter 平台上线，
被描述为“frontier omni-modal model”，
支持视觉、听觉和推理的统一处理。
其身份尚未公开确认，但其能力表现
引发了关于多模态模型评估标准化的广泛讨论。

\subsection{从“外挂视觉”到“原生多模态”：架构范式的根本转变}

2025--2026 年多模态领域最深刻的技术转变，
不是某个基准的刷新，
而是\textbf{架构范式从“外挂”到“原生”的迁移}。
理解这一转变，需要追溯多模态模型的三代架构演进。

\textbf{第一代：转录-理解管线}（2022--2023）。
视觉和语言是两个完全独立的系统。
图像先由 CLIP/ViT 编码为向量，
语音先由 Whisper 转录为文本，
然后这些“翻译后的”信号被送入 LLM。
这种架构的本质是\textbf{模态翻译}：
将非文本模态翻译成文本能理解的形式。
信息损失巨大：图像中的空间关系被压缩为一维向量，
语音中的韵律和情感在转录过程中完全丢失。

\textbf{第二代：投影融合}（2023--2024）。
LLaVA~\cite{liu2023llava} 范式的核心创新是
用一个可训练的投影层将视觉 token 映射到 LLM 的输入空间，
然后与文本 token \textbf{拼接}送入同一个 Transformer。
这比第一代进步显著，
视觉信息以更细粒度的 token 形式被 LLM 处理，
而不是压缩为单一向量。
但架构上的根本限制是：
\textbf{LLM 是在纯文本上预训练的}。
视觉编码器在 LLM 预训练完成后才被“接入”，
通过投影层和有限的多模态微调数据来“学会”
如何处理视觉 token。
这就像一个只读过书的人第一次看到图片，
他可以描述图片中的物体，
但缺少从出生起就积累的视觉直觉。

\textbf{第三代：原生多模态}（2024--2026）。
GPT-4o、Gemini 3（基于 Gemini 架构~\cite{anil2023gemini}）
和 Claude 4.5 Sonnet 代表了根本性的范式转变——
\textbf{模型从预训练的第一天起就同时处理文本、图像和音频}。
这不是工程细节的差异，
而是模型对世界的理解方式的根本不同。

从架构角度看，原生多模态模型的关键设计选择包括：

\textbf{统一的 tokenizer}。
所有模态的输入被映射到\textbf{同一个}离散 token 空间中。
图像通过学习的视觉 tokenizer（如 VQ-VAE 或 dVAE）
被离散化为视觉 token，
音频通过声学 tokenizer（如 SoundStream 或 EnCodec）
被离散化为音频 token，
然后这些 token 与文本 token
在同一个 Transformer 中交织处理。
Gemini 3 的架构~\cite{anil2023gemini}在此方向上走得最远，
它在预训练阶段就将文本、图像、音频和视频数据
混合在同一个训练批次中，
使模型学到了\textbf{跨模态的深层语义关联}
（例如，“玻璃碎裂的声音”与“碎玻璃的图像”
在模型的内部表示空间中是近邻）。

\textbf{交叉注意力 vs.\ 自注意力}。
在处理多模态输入时，
有两种主要的注意力设计。
交叉注意力（cross-attention）方案
让不同模态的 token 通过专门的交叉注意力层交互，
文本 token 可以“查询”视觉 token 以获取视觉信息。
Flamingo~\cite{alayrac2022flamingo} 是这种设计的代表。
自注意力（self-attention）方案
则将所有模态的 token 简单拼接，
通过标准的自注意力让它们自由交互，
GPT-4o 和 Gemini 采用了这种更简洁但计算成本更高的设计。
自注意力方案的优势是
\textbf{模型可以学到任意模态之间的关系}——
不受预定义的交互模式限制。

\textbf{原生多模态的训练成本}是一个关键考量。
与纯文本预训练相比，
多模态预训练需要数量级更多的数据
（互联网上的文本远多于高质量的图文配对数据），
更长的训练时间
（视觉 token 的数量通常远大于对应的文本描述），
以及更复杂的数据管道
（图像/视频/音频数据需要不同的预处理和增强策略）。
据估算，GPT-4o 的多模态预训练成本
是纯文本 GPT-4 预训练成本的 3--5 倍。
这也是为什么原生多模态模型
目前仍主要由资源最充裕的少数实验室
（OpenAI、Google、Anthropic、字节跳动）掌握。

\textbf{Claude 的多模态路线}值得单独讨论。
Anthropic 在多模态方面采取了更为审慎的策略，
Claude 3 系列（2024 年 3 月）首次引入视觉能力，
Claude 3.5 Sonnet 将视觉理解提升到 SOTA 水平，
Claude 4.5 Sonnet（2025 年 2 月）
进一步成为多项视觉基准的最佳模型。
然而，Claude 系列至今没有引入原生图像生成能力，
也没有引入原生音频输入/输出，
Anthropic 选择将工程资源集中在
\textbf{将每种已支持的模态做到极致}——
而非追求全模态覆盖。
这种策略在 Agent 场景中有其合理性，
Claude Computer Use 需要的是
对屏幕截图的极精确视觉理解，
而不是生成图像或处理音频的能力。

\begin{whybox}[原生多模态为什么重要？]
一个在预训练阶段就同时处理视觉和语言的模型，
与一个先学会语言再“学看图”的模型，
差异不仅在于基准分数。
关键差异在于\textbf{接地性}（grounding），
原生多模态模型将语言概念直接锚定在视觉经验上。
当它处理“fragile”（易碎的）这个词时，
它的内部表示不仅编码了词典定义，
还融合了无数个玻璃碎裂、陶瓷摔碎的视觉片段。
这种接地性使模型在理解涉及物理世界的语言描述时
更加准确和可靠。
人类的语言理解深度依赖于感官经验：
一个从未见过“红色”的人很难真正理解“红色”这个词。
原生多模态模型正在弥补这种
纯文本模型天生缺乏的感官基础。
\end{whybox}


\section{AI for Science：LLM 技术向科学发现的延伸}
\label{sec:ai-for-science}

LLM 背后的核心技术，Transformer 架构、
大规模预训练、自回归生成，
正在被迁移到远超自然语言处理的科学领域。
这不是简单的“用 GPT 回答科学问题”，
而是用类似的方法论\textbf{从头训练科学领域的基础模型}——
直接在分子结构、蛋白质序列、天气数据和材料属性上
学习自然界的规律。

\subsection{蛋白质结构预测与设计}

\textbf{AlphaFold}~\cite{jumper2021alphafold} 的故事
是 AI for Science 最具标志性的成功案例。
2020 年，AlphaFold 2 在 CASP14 竞赛中
以远超所有竞争对手的精度预测蛋白质三维结构，
解决了一个困扰生物学界 50 年的难题。
2023 年的 AlphaFold 3~\cite{abramson2024alphafold3}
进一步将预测范围从单体蛋白质
扩展到蛋白质-蛋白质相互作用、
蛋白质-DNA/RNA 复合物以及蛋白质-小分子结合——
这对药物发现至关重要，
因为大多数药物通过与特定蛋白质结合来发挥作用。

AlphaFold 3 的架构变化值得注意：
它从 AlphaFold 2 的\textbf{迭代精化}（Evoformer + Structure Module）
转向了\textbf{扩散模型}——
使用扩散过程从噪声中逐步“去噪”出原子坐标。
这与 Sora 等视频生成模型使用的扩散范式异曲同工，
都是从随机噪声开始，
通过学习的去噪过程逐步生成目标结构。
这种架构迁移暗示了一个更深层的统一：
\textbf{无论是像素、蛋白质还是分子，
“从噪声中生成结构”都是一种普适的生成范式}。

在蛋白质设计方面（“逆向”问题，给定功能，设计结构），
\textbf{RFdiffusion}~\cite{watson2023rfdiffusion}（David Baker 实验室，2023 年）
将扩散模型用于从头设计蛋白质骨架，
能够生成满足特定几何约束的全新蛋白质结构。
\textbf{ProteinMPNN}~\cite{dauparas2022proteinmpnn}
则解决了“逆折叠”问题，
给定一个蛋白质骨架结构，设计最优的氨基酸序列。
这两个工具的组合使得
从“我想要一个能结合 X 蛋白质的分子”
到“这是具体的氨基酸序列”的全链路自动化成为可能。

2024 年 10 月，David Baker 和 Demis Hassabis
因“蛋白质结构预测和蛋白质设计”的贡献
共同获得诺贝尔化学奖，
这是 AI 技术首次直接促成诺贝尔奖级别的科学发现。

\subsection{材料科学与分子发现}

\textbf{GNoME（Graph Networks for Materials Exploration）}%
~\cite{merchant2023gnome}（Google DeepMind，2023 年 11 月）
使用图神经网络发现了 220 万种新的稳定晶体结构，
相当于人类在此前 800 年的实验中发现数量的 45 倍。
GNoME 的训练数据来自 Materials Project 等开源材料数据库，
通过学习已知稳定材料的原子排列模式，
预测哪些尚未被合成的材料组合可能是热力学稳定的。
其中 381,000 种材料已经通过后续实验验证~\cite{merchant2023gnome}，
部分材料展现出了用于下一代电池、太阳能电池
和超导体的潜在价值。

分子生成模型
（用于药物先导化合物发现和材料设计）
也在快速发展。
\textbf{Uni-Mol 2}（DP Technology，2024 年）
将 Transformer 用于 3D 分子表示学习，
在分子属性预测和 3D 构象生成上达到了 SOTA。
AlphaFold 3 的扩散架构启发了
一系列新的分子生成模型，
直接在三维空间中“生长”分子结构，
同时满足化学键约束和药物性质要求。

\subsection{天气与气候预测}

天气预测是 AI for Science 的另一个突破性领域。
传统的数值天气预报（NWP）
通过求解大气的偏微分方程组来预测天气变化，
需要超级计算机运行数小时。
AI 天气模型正在以秒级推理时间
达到甚至超越传统 NWP 的精度。

\textbf{GenCast}~\cite{price2024gencast}（Google DeepMind，2024 年 12 月）
使用扩散模型生成集合天气预报，
不是给出一个确定的预测，
而是生成一组可能的天气演变路径，
自然地量化了预测的不确定性。
GenCast 在 97.2\% 的评估指标上
超越了顶级传统 NWP 系统 ENS
（欧洲中期天气预报中心的集合预报系统），
并且推理速度快约 1000 倍：
生成 15 天的全球天气预报仅需 8 分钟
（在单块 TPU 上），
而 ENS 需要数小时的超级计算机运算。

\textbf{Pangu-Weather}~\cite{bi2023pangu}（华为，2023 年）
是首个在精度上全面超越传统 NWP 的 AI 天气模型，
使用 3D Earth-Specific Transformer 架构
直接在地球的经纬度网格上建模大气状态演变。
后续的 \textbf{FuXi}（复旦大学）和
\textbf{FourCastNet}（NVIDIA）
进一步拓展了 AI 天气预测的能力边界。

值得注意的是，这些 AI 天气模型使用的架构
本质上都是 Transformer 或其变体，
它们将大气状态表示为空间-时间 token，
通过注意力机制学习大气动力学的模式。
训练数据则来自 ERA5 再分析数据集
（欧洲气象中心基于历史观测数据重建的全球大气状态数据）。
\textbf{一个在视频上训练的 Transformer
和一个在天气数据上训练的 Transformer，
在架构上几乎没有区别}，
差异在于输入的“帧”不是 RGB 像素，
而是温度、气压、风速等物理量。

\subsection{AI 辅助的药物发现}

药物发现是 AI 最具商业价值的科学应用之一。
2025 年，首批\textbf{完全由 AI 设计的药物}
进入了临床试验阶段。
\textbf{Insilico Medicine} 的 INS018\_055
（治疗特发性肺纤维化的小分子药物）
是全球首个从靶点发现到化合物设计全程由 AI 主导的
进入二期临床试验的候选药物。
从靶点识别到候选分子提名仅耗时 18 个月，
传统流程通常需要 4--5 年~\cite{zhavoronkov2019deep}。

AI 在药物发现中的作用不限于分子设计。
\textbf{靶点发现}——
通过分析基因组数据和疾病关联，
AI 识别出传统方法难以发现的治疗靶点。
\textbf{ADMET 预测}（吸收、分布、代谢、排泄、毒性）：
在分子被合成之前就预测其在人体内的行为，
大幅减少昂贵的实验筛选。
\textbf{临床试验设计}——
AI 优化患者招募策略和剂量方案，
提高试验成功率。

然而，AI 药物发现也面临独特的挑战：
\textbf{数据稀缺}——
临床数据量远小于互联网文本数据。
整个 ChEMBL 数据库（最大的生物活性分子数据库之一）
包含约 200 万个化合物和 2,000 万个活性数据点，
对于 LLM 的标准来说，这是微不足道的数据量。
相比之下，GPT-4 的预训练语料包含超过 10 万亿 token。
这种数据规模的差异意味着
AI 药物发现模型必须在\textbf{远少于 LLM}的数据量上
学会化学和生物学的规律，
这对模型的归纳偏置（inductive bias）设计
和数据效率提出了更高的要求。

\textbf{分布外泛化}——
模型需要生成训练数据中从未出现的分子，
而不仅仅是组合已知的模式。
这与 LLM 的“创造性写作”在本质上不同：
一个语法正确但语义荒谬的句子可以被读者轻松识别，
但一个化学上合理但药理学上无效的分子
只能通过昂贵的实验来验证。
“在化学空间中创新”比“在语言空间中创新”
需要更严格的约束和更可靠的评估。

\textbf{端到端验证}：
从计算预测到实验验证到临床试验的链条极长，
反馈循环以年为单位而非以秒为单位。
一个 AI 生成的候选分子从计算设计到完成二期临床试验
通常需要 5--7 年。
即使 Insilico Medicine 将发现阶段从 4--5 年压缩到 18 个月——
后续的临床试验仍然受制于生物学的时间尺度。

\textbf{Recursion Pharmaceuticals}（2025 年完成与 Exscientia 的合并）
代表了 AI 药物发现的另一种路线，
\textbf{实验室自动化 + AI 闭环}。
Recursion 运营着高度自动化的生物实验室，
每周可以生成数百万张细胞显微图像，
这些图像被 AI 模型分析以发现药物-疾病的关联。
这种“让 AI 设计实验，让机器人执行实验，
让 AI 分析结果并设计下一轮实验”的闭环
是 AI 在实验科学中最完整的应用范式。

\subsection{AI for Science 的基础设施挑战}

科学 AI 面临的基础设施挑战
与 LLM 既有相似也有独特之处。

\textbf{科学数据的异构性}。
与互联网文本的相对标准化不同，
科学数据高度异构，
不同实验室使用不同的测量设备、
不同的数据格式和不同的标注标准。
将分散在全球数万个实验室的科学数据
整合为统一的训练集，
其工程难度堪比互联网时代之前的“搜索引擎问题”。

\textbf{物理约束的编码}。
科学模型必须遵守物理定律：
生成的蛋白质结构必须满足化学键角约束，
预测的天气状态必须满足质量和能量守恒。
如何将这些“硬约束”嵌入到
本质上是概率性的神经网络中，
是一个活跃的研究方向。
当前的主流方法包括：
在损失函数中添加物理约束惩罚项、
在生成后进行物理修正
（post-hoc correction）、
以及设计具有内置对称性的网络架构
（如等变神经网络 / equivariant neural networks）。

\textbf{可解释性}。
科学家不仅需要预测结果，
还需要\textbf{理解为什么}。
一个能够准确预测蛋白质结构但无法解释
“为什么这个氨基酸序列折叠成这个形状”的模型，
对科学理解的贡献是有限的。
这与 LLM 领域的“黑箱”问题一脉相承，
但在科学场景中更为紧迫，
因为科学进步的核心不是预测，而是理解。

\begin{keybox}[AI for Science 的共同模式]
回顾上述各领域，一个清晰的模式浮现：
（1）\textbf{将科学数据 tokenize}——
蛋白质序列用氨基酸 token 表示，
分子用 SMILES 字符串或 3D 坐标表示，
天气用经纬度网格上的物理量 token 表示；
（2）\textbf{用 Transformer 学习 token 之间的关系}——
无论是蛋白质中氨基酸的空间相互作用、
分子中原子的化学键、
还是大气中不同位置的气象关联；
（3）\textbf{用生成模型（自回归或扩散）创造新结构}：
生成新的蛋白质、分子或天气预报。
这个“tokenize → learn → generate”的三步范式
是 LLM 技术向科学延伸的通用框架。
LLM 训练的核心洞察——
“用足够的数据和计算，
让复杂性从简单的预测目标中涌现”：
在科学领域同样适用。
\end{keybox}


\section{Agent：从对话到行动}

LLM 正在从「回答问题」演进为「执行任务」。
这是 2025--2026 年最重要的技术趋势之一。

\subsection{Agent 的核心架构}

Agent 系统的核心是一个循环：
\textbf{观察}（Observe）→
\textbf{推理}（Think）→
\textbf{行动}（Act）→
\textbf{观察}（Observe）。
这就是 ReAct 范式——
模型交替进行推理（生成文本思考过程）
和行动（调用工具获取外部信息），
直到任务完成或达到预设的步数上限。

这个循环可以持续数十到数百步。
对于复杂任务（如修复一个涉及多文件的 GitHub issue），
Agent 可能需要：
阅读代码 → 理解架构 → 定位 bug →
编写修复 → 运行测试 → 分析失败 → 修改修复 → 再次测试。
每一步都是「观察-思考-行动」的实例。

Agent 架构的另一个关键组件是\textbf{记忆系统}。
短期记忆（对话上下文）受限于模型的上下文窗口；
长期记忆（跨会话的知识积累）
则需要外部存储机制，如向量数据库、
结构化知识图谱或简单的文件系统。
2025--2026 年的趋势是将记忆系统\textbf{层次化}：
工作记忆（当前任务的中间结果）、
情景记忆（过去任务的经验和教训）、
语义记忆（领域知识和用户偏好）：
每一层使用不同的存储和检索策略。
这种层次化设计使 Agent 能够
在长时间跨度内保持一致的行为和知识积累。

\subsection{工具使用与标准化}

Agent 的能力取决于它可以使用的工具。
当前的主流工具类型包括：

\begin{itemize}[noitemsep]
  \item \textbf{Function calling}：
        LLM 生成结构化的函数调用
        （指定函数名和参数），
        由外部系统执行并返回结果。
        这是最基础的工具使用形式。
  \item \textbf{Code execution}：
        LLM 生成代码（Python、Bash 等），
        在沙箱环境中执行并获取输出。
        Code Interpreter（OpenAI）和 tool use（Anthropic）
        是典型实现。
  \item \textbf{Web browsing}：
        LLM 控制浏览器访问网页，提取信息。
        需要处理动态渲染、验证码、反爬虫等挑战。
  \item \textbf{Computer use}：
        LLM 通过理解屏幕截图和控制鼠标/键盘
        来操作计算机。
        Anthropic 的 Claude computer use 是这个方向的先驱，
        模型可以打开应用程序、填写表单、点击按钮，
        像人类用户一样操作任意桌面软件。
\end{itemize}

\textbf{MCP（Model Context Protocol）}
由 Anthropic 于 2024 年底开发并开源，
作为 Agent 与外部工具交互的标准协议，
正在被越来越多的工具提供商采用。
它定义了工具描述（名称、参数、返回值）、
调用和结果返回的标准格式，
使得同一个 Agent 可以无缝对接不同的工具和服务，
搜索引擎、数据库、API、本地文件系统、
甚至其他 AI 模型。
MCP 的意义类似于 HTTP 之于 Web：
一个统一的协议使得生态系统可以繁荣发展。

\textbf{MCP 生态系统的爆发式增长}（2025 年底--2026 年初）
验证了这一判断。
截至 2026 年 3 月，MCP 相关工具的下载量已超过 9700 万次。
2025 年 12 月，Anthropic 将 MCP 捐赠给
Linux 基金会旗下的 Agentic AI Foundation，
使其成为由行业共同治理的中立标准。
2026 年 1 月，OpenAI、Google 和 Microsoft
先后宣布在各自的 Agent 框架中采用 MCP，
这标志着 MCP 从 Anthropic 的专有协议
正式成为行业事实标准（de facto standard）。
2026 年 1 月 26 日发布的 \textbf{MCP Apps}
进一步扩展了 MCP 的能力边界，
Agent 不再只能调用后端 API，
还可以通过 MCP Apps 呈现交互式 UI 组件，
使 Agent 的输出从纯文本升级为富交互体验。

\subsection{代码 Agent：最成熟的落地场景}
\label{sec:code-agents}

2025--2026 年，Agent 应用最成熟的领域是\textbf{软件工程}。
代码 Agent 已经从辅助工具演进为半自主的编程助手，
在某些场景下甚至可以独立完成完整的开发任务。

\textbf{Claude Code}（Anthropic，2025 年 2 月发布）
是 CLI 级别的自主编程 Agent。
与传统的 IDE 插件不同，
Claude Code 运行在终端中，
可以直接读写文件、执行 bash 命令、运行测试、
管理 git 操作，全部通过自然语言交互。
它支持 “headless mode”（无人值守模式），
可以在 CI/CD 管道中自动执行代码审查、
修复 lint 错误、生成测试等任务。
Claude Code 的设计理念是
\textbf{给予 Agent 与人类开发者相同的工具和权限}——
不是通过 API 读取代码片段，
而是像人类一样在真实的开发环境中工作。
\textbf{Claude Code 2.0}（2026 年 2 月 24 日）
将这一理念推向了新阶段：
引入多 Agent 编排（orchestration），
主 Agent 可以将子任务分派给专门化的子 Agent；
持久化项目记忆（persistent project memory），
跨会话保留项目上下文和决策历史；
原生 CI/CD 集成，
Agent 可以直接触发构建、测试和部署流水线，
将代码 Agent 从辅助工具升级为
软件工程工作流的核心参与者。

\textbf{Devin}（Cognition AI）
是另一种路线的代码 Agent，
完全自主的「AI 软件工程师」。
Devin 在云端独立运行，
用户通过 Slack 或 Web 界面分配任务，
Devin 自主规划、编码、测试，
最终以 Pull Request 的形式交付结果。
Devin 适合被委派独立的、
定义清晰的开发任务（如 “将这个 API 的认证从 OAuth 1.0 迁移到 OAuth 2.0”）。
2025 年底 Cognition 收购了 Windsurf IDE，
将 Devin 的自主能力与 IDE 的交互式编程结合。

\textbf{Cursor} 和 \textbf{Windsurf}（Codeium）
代表了 IDE-integrated Agent 的路线：
将 AI 能力深度嵌入代码编辑器。
Cursor 的 “Composer” 模式可以跨多文件进行修改，
Agent 模式可以自动运行终端命令和管理文件。

SWE-bench（修复真实 GitHub issue）
的解决率已经从 2024 年初的 $\sim$5\%
提升到 2026 年初的 $\sim$79\%
（Claude Opus 4.6 Thinking 在 SWE-bench Verified 上的成绩），
MiniMax M2.5 更是达到了 80.2\%。
然而，SWE-bench Verified 的高分可能部分受到数据污染的影响。
更能反映真实能力的 \textbf{SWE-bench Pro}
（包含 1,865 个多语言问题的更严格版本）
显示当前最优模型的解决率仅为 $\sim$44\%，
这表明从 benchmark 到真实软件工程仍有相当距离。

\begin{corebox}[代码 Agent 的使用范式演变]
代码 Agent 的使用范式正在经历三个阶段：
\textbf{阶段一：补全}——GitHub Copilot 风格，
逐行或逐函数补全代码，人类主导、AI 辅助。
\textbf{阶段二：对话式编程}——
通过自然语言描述需求，Agent 生成或修改代码块，
人类审查并接受/拒绝。Cursor、Claude Code 当前主要在这个阶段。
\textbf{阶段三：委托式编程}：
将完整的任务委托给 Agent，
Agent 自主完成规划、实现和测试，
人类仅在关键决策点介入。
Devin 的设计理念指向这个阶段，
但可靠性仍然是核心挑战。
从阶段二到阶段三的跨越
需要 Agent 具备更强的规划能力、
错误恢复能力和对不确定性的处理能力。
\end{corebox}

\subsection{多 Agent 系统}

更复杂的任务可以由多个专门化的 Agent 协作完成。
例如，一个软件开发任务可以分配给：
\textbf{架构师 Agent}（设计整体方案）、
\textbf{开发者 Agent}（编写代码）、
\textbf{测试 Agent}（编写和运行测试）、
\textbf{审查 Agent}（代码审查和质量检查）。

多 Agent 框架包括
AutoGen（Microsoft）、CrewAI、LangGraph 等。
它们提供了 Agent 间通信、任务分配、
结果汇聚等基础设施。

多 Agent 系统的核心挑战是\textbf{协调成本}——
Agent 之间的通信（通常以自然语言进行）
本身消耗大量 tokens，
而且 Agent 间的误解可能导致任务偏离方向。
早期的研究表明，
对于大多数任务，一个能力足够强的单 Agent
比多个能力较弱的 Agent 协作更有效。

然而，2026 年初的进展正在改变这一判断。
\textbf{Kimi K2.5 Agent Swarm}（2026 年 1 月 27 日）
是多 Agent 协作的一个重要突破，
通过 PARL（Parallel Agent Reinforcement Learning）协议，
一个主 Agent 可以协调多达 100 个子 Agent 并行执行任务，
单次任务最长可运行 1500 步。
关键创新在于 PARL 协议减少了 Agent 间的自然语言通信开销，
转而使用结构化的任务分解和结果汇聚机制，
实现了 80\% 的运行时间缩减。
这表明多 Agent 系统的效率瓶颈不在于概念本身，
而在于协调协议的设计。

\textbf{OpenRouter Hunter Alpha}（2026 年 3 月 11 日）
是另一个值得关注的信号，
这个以匿名形式上线的模型被描述为
“1T 参数 + 1M 上下文”的 Agent 专用模型，
其规格与 DeepSeek V4 的传闻高度吻合。
无论其真实身份如何，
“为 Agent 而生”的模型设计理念
正在成为下一代基础模型的核心考量。

\subsection{Agent 的可靠性挑战}

Agent 面临的最大挑战是\textbf{可靠性}。
在多步骤任务中，每一步都有一定的失败概率。
如果每步的成功率是 95\%，
那么 20 步链式执行的整体成功率只有 $0.95^{20} \approx 36\%$。

更危险的是\textbf{静默失败}——
Agent 可能在某一步产生了幻觉
（比如调用了一个不存在的 API，
或者基于错误的假设做出了决策），
但后续步骤继续在错误的基础上执行，
最终产生看似合理但完全错误的结果。

幻觉在工具使用中尤其危险：
在对话中，用户可以识别不合理的回答；
但在 Agent 自主执行的过程中，
一个错误的数据库查询或文件操作可能造成真实的损害。

WebArena（自主网页操作）和 OSWorld（操作系统级任务）
等基准测试正在推动 Agent 能力的边界。
但从 benchmark 到真实世界仍有差距，
benchmark 中的任务有明确的开始和结束条件，
而真实世界的任务往往目标模糊、
需要人类反馈、容错要求更高。

\textbf{可靠性差距}（Reliability Gap）
是 Agent 商业化的核心挑战。
在 benchmark 上 80\% 的成功率意味着
5 次执行中有 1 次失败，
对于学术评估这是优秀的成绩，
但对于生产环境这意味着
\textbf{每 5 个自动化任务中就有 1 个需要人工干预}。
许多企业场景要求 95\%+ 的可靠性
才能证明自动化的经济价值，
因为人工干预本身有成本，
如果干预频率过高，
自动化带来的效率提升会被干预成本抵消。

应对可靠性差距的工程策略包括：
\textbf{多轮重试}——
Agent 在检测到失败后自动重试，
利用 LLM 的非确定性（相同输入的不同采样路径）
获得不同的执行策略；
\textbf{自验证}——
Agent 在完成任务后自行验证结果
（如运行测试、检查输出格式），
不满足验证条件则回退到上一个检查点；
\textbf{人在环路的渐进式撤退}——
Agent 从需要频繁人类确认的模式开始，
随着信任的建立逐步减少确认频率。
Claude Code 的权限系统
（\texttt{allow}/\texttt{deny}/\texttt{ask} 规则）
正是这种渐进式信任模型的工程实现。

\begin{warnbox}[Agent 的经济影响]
Agent 作为「数字劳动力」（digital workers）的概念
正在引发关于工作未来的严肃讨论。
如果一个 Agent 能够完成初级软件工程师 50\% 的任务，
这对就业市场意味着什么？
乐观者认为 Agent 会放大人类的生产力，
创造新的工作类型（如 Agent 监督者、prompt 工程师）。
悲观者担心大规模的工作岗位替代。
现实可能介于两者之间——
Agent 的可靠性问题限制了它的自主性，
但它正在显著改变工作的内容和效率要求。
Anthropic CEO Dario Amodei 在 2025 年的文章
“Machines of Loving Grace” 中
勾勒了 AI 增强人类而非替代人类的愿景，
但同时承认转型期的阵痛不可避免。
\end{warnbox}

\section{具身智能：从屏幕走向物理世界}
\label{sec:embodied-ai}

LLM 技术的下一个前沿不在数字世界，而在\textbf{物理世界}。
具身智能（Embodied AI）：
让 AI 系统在机器人、无人机、自动驾驶等物理载体上
感知环境、做出决策并执行动作——
正在从学术研究走向工程落地。

\subsection{大语言模型作为机器人的“大脑”}

传统的机器人控制依赖于手工编写的规则和状态机，
\texttt{if sensor\_reading > threshold then execute\_action}。
这种方法在工厂流水线等结构化环境中有效，
但面对开放世界的复杂性和不确定性时极其脆弱：
一个没有被预编程应对的场景就可能导致系统完全卡死。

LLM 为机器人提供了一种全新的能力：
\textbf{将自然语言指令转化为动作序列}。
用户说“帮我把桌上的红色杯子放到厨房水槽里”，
LLM 可以将这条指令分解为
一系列可执行的子任务，
定位桌子 → 识别红色杯子 → 抓取 → 导航到厨房 →
定位水槽 → 放置。
这种从语言到动作的“翻译”能力
是 LLM 在具身智能领域最直接的贡献。

\textbf{RT-2（Robotic Transformer 2）}~\cite{brohan2023rt2}
（Google DeepMind，2023 年）
是这一思路的里程碑式工作。
RT-2 直接在 vision-language 模型（PaLM-E）上
微调机器人操作数据，
模型的“输出 token”不再只是文本，
还包括机器人关节角度和末端执行器位置的\textbf{离散化动作 token}。
这意味着\textbf{同一个 Transformer
既能回答“这是什么物体？”
也能输出“将机械臂移动到 (x, y, z) 坐标”的动作指令}。
RT-2 展现了令人惊讶的泛化能力，
它能够执行训练数据中从未出现过的物体操作，
因为 vision-language 模型的通用知识
（“瓶子可以用来浇水”）
被迁移到了机器人操作场景中。

后续的 \textbf{Octo}（UC Berkeley，2024 年）
和 \textbf{OpenVLA}（Stanford/TRI，2024 年）
进一步推进了机器人基础模型的方向，
它们在来自多种机器人平台、多种任务的
大规模操作数据集上预训练，
然后通过少量任务特定数据微调到具体的机器人和任务。
这与 LLM 的“预训练 + 微调”范式完全一致。

\subsection{人形机器人：硬件载体的突破}

2025--2026 年，人形机器人从实验室走向了商业化的早期阶段。

\textbf{Figure 02}（Figure AI，2025 年发布）
是将 LLM 与人形机器人深度整合的代表作。
Figure 02 搭载了 OpenAI 提供的视觉-语言模型，
能够通过自然语言与人类对话，
理解复杂的操作指令，
并在物理世界中执行多步骤任务。
Figure AI 在 2024 年完成了 6.75 亿美元融资
（投资方包括 Microsoft、NVIDIA、OpenAI、Jeff Bezos）——
估值 26 亿美元，
这表明资本市场对“LLM + 人形机器人”路线的信心。
Figure 与 BMW 等制造商合作，
目标是在汽车装配线等场景中部署人形机器人。

\textbf{Tesla Optimus}（Tesla，2022 年原型亮相，
2025 年进入有限试生产）
是消费级人形机器人的最大胆尝试。
Elon Musk 声称 Optimus 的长期目标是
“每个人都能拥有的通用机器人助手”，
目标售价低于 2 万美元。
截至 2026 年初，Optimus 已在 Tesla 工厂内部
执行简单的物料搬运和零件分拣任务，
但距离“通用家庭助手”的愿景仍有巨大差距。
Tesla 的独特优势在于其\textbf{自动驾驶数据管道}——
FSD（Full Self-Driving）积累的
数十亿英里真实世界行驶数据和训练基础设施
可以被迁移到人形机器人的训练中。

\textbf{1X Neo}（1X Technologies，2025 年发布原型）
是挪威公司 1X Technologies 的人形机器人，
专注于家庭和商业服务场景。
Neo 的设计哲学与 Figure/Tesla 不同，
它更强调\textbf{安全的人机共存}
而非纯粹的力量和速度，
使用柔性驱动器（compliant actuators）
确保在意外碰撞时不会伤害人类。
OpenAI 是 1X 的主要投资方——
这暗示了 OpenAI 在具身智能领域的战略布局。

\textbf{Unitree H1 / G1}（宇树科技，2024--2025 年）
是中国人形机器人的代表，
以“高性价比 + 运动能力”著称。
Unitree H1 在公开视频中展示了
全地形行走、跑步、跳跃甚至后空翻的能力，
其硬件成本据称远低于西方竞品，
这与中国在消费电子制造方面的传统优势一致。
Unitree 的开源控制框架
吸引了大量学术研究者使用其平台进行研究。

\subsection{Sim-to-Real：模拟到现实的迁移}

在物理世界中训练机器人是昂贵且缓慢的：
每次失败都可能损坏硬件，
数据收集速度受限于物理时间。
\textbf{Sim-to-Real Transfer}（从模拟到现实的策略迁移）
是解决这一问题的主要方法。

\textbf{NVIDIA Isaac Sim / Isaac Lab}~\cite{makoviychuk2021isaac}
是当前最广泛使用的机器人模拟平台。
Isaac Sim 基于 NVIDIA Omniverse 构建，
提供了光线追踪级别的视觉真实性
和 PhysX 引擎驱动的物理模拟。
关键的训练范式是 \textbf{Domain Randomization}：
在模拟中随机化物体的颜色、纹理、大小、
光照条件和物理参数
（摩擦系数、重力、传感器噪声），
迫使策略网络学习对这些变化\textbf{鲁棒}的行为，
从而在迁移到真实世界时不会因“看起来不一样”而失效。

Isaac Gym（Isaac Lab 的前身）
在强化学习训练中实现了惊人的并行性：
在单块 GPU 上同时运行数千个模拟环境，
一天内完成相当于数年真实世界经验的训练。
这种\textbf{计算换时间}的策略
与 LLM 训练中的 “scaling through compute” 思路一致。

Sim-to-Real 的核心挑战是 \textbf{Reality Gap}——
模拟环境与真实世界之间不可避免的差异。
无论模拟器多么精细，
它都无法完全捕捉真实世界的所有物理细节，
接触力学、柔性材料变形、
液体交互、传感器噪声的真实分布。
当前的缓解策略包括：
更精确的物理引擎、
更大范围的 Domain Randomization、
以及“Sim-to-Real-to-Sim”循环
（在真实世界收集少量数据来校准模拟器参数）。

\begin{corebox}[具身智能的“GPT 时刻”还有多远？]
具身智能领域正在寻找自己的“GPT 时刻”：
一个像 ChatGPT 一样彻底改变公众认知的突破。
目前的人形机器人仍然在以下方面存在根本限制：
（1）\textbf{灵巧操作}——
精细的手部操作（如折叠衣物、系鞋带）
远比行走和搬运困难；
（2）\textbf{长时间自主性}——
电池续航和计算资源限制了
机器人的持续工作时间；
（3）\textbf{安全性}：
在人类身边工作的机器人
必须在所有可能的故障模式下保证安全。
乐观的估计是 2027--2028 年
人形机器人将在特定场景（如仓库物流、制造业）
实现规模化商业部署。
但“通用家庭机器人”，能够在任意家庭环境中
执行任意日常任务，
可能仍需 10 年以上的技术积累。
\end{corebox}


\section{硬件路线图：算力的下一个台阶}
\label{sec:hardware-roadmap}

LLM 的能力演进与硬件进步密不可分。
2025--2027 年的 AI 加速器路线图
决定了下一代模型的能力上限。

\subsection{NVIDIA：从 Blackwell 到 Vera Rubin}

NVIDIA 继续以 1--2 年的节奏推出新一代 GPU 架构：

\textbf{Blackwell 架构}（2024 年发布，2025 年量产）
是当前一代的旗舰。
B200 GPU 提供 9 PFLOPS 的 FP4 性能
和 192GB HBM3e 显存（6.4 TB/s 带宽），
相比上一代 H100（FP8 约 3.9 PFLOPS，80GB HBM3）
有显著提升。
\textbf{GB200 NVL72} 是 Blackwell 的系统级产品：
将 72 块 B200 GPU 通过第五代 NVLink 互连
（GPU 间带宽 1.8 TB/s）组成一个液冷机柜，
提供 720 PFLOPS 的 FP4 计算能力和 13.5TB 的统一显存。
这种设计的意义在于：
一个 GB200 NVL72 机柜可以将一个 700B+ 参数的模型
完全放入 GPU 显存中进行推理，
消除了多机通信的延迟。

\textbf{Vera Rubin 架构}（2026 年 1 月 CES 正式发布）
是 NVIDIA 的下一代平台。
Rubin GPU 采用 3nm 工艺（2025 年 8 月已 tape out），
搭配 HBM4 内存（每 GPU 可能达到 256GB+，
带宽超过 8 TB/s）。
NVIDIA 声称 Vera Rubin 将实现推理成本降低至 Blackwell 的 1/10，
性能提升 5 倍，这意味着推理经济学将发生根本性变化。
Vera Rubin Ultra 将是完整平台版本，
包含 CPU（Vera，基于 ARM 架构）和 GPU（Rubin）的集成。

\subsection{AMD：从追赶者到竞争者}

AMD 的 Instinct 系列正在成为 NVIDIA 的真正竞争对手：

\textbf{MI300X}（2024 年量产）
提供 192GB HBM3 显存和 5.3 TB/s 带宽：
在显存容量上超越了 H100。
这使得 MI300X 在大模型推理场景中有显著优势
（更大的模型可以放入单 GPU），
微软 Azure 和 Oracle Cloud 都已部署 MI300X 集群。

\textbf{MI350 系列}（2025 年底--2026 年初量产）
基于全新的 CDNA 4 架构和 3nm 工艺，
采用 chiplet 设计。
旗舰产品 \textbf{MI355X DLC} 以液冷机架形式交付，
单机架 128 块 GPU，提供 36TB HBM3e 显存
和 2.6 exaflops 的 FP8 计算能力，
直接对标 NVIDIA GB200 NVL72。
AMD 在 Hot Chips 2025 上展示了 MI350 的 node-to-rack 扩展架构，
MI350 的关键卖点是
\textbf{开源软件栈}——ROCm 生态的持续改善
使得越来越多的训练框架
（PyTorch、JAX、vLLM）原生支持 AMD GPU，
降低了迁移成本。

\textbf{MI400}（路线图上为 2026 年底--2027 年）
将是 AMD 的下一个大跳跃，
基于 2nm 工艺，目标是达到 yottaflops 级别的系统计算能力，
这将是 AI 加速器领域有史以来最激进的性能目标。

\subsection{Google TPU：垂直整合的典范}

Google 的 TPU（Tensor Processing Unit）
代表了另一种路线——
自研加速器完全为自己的工作负载优化。

\textbf{TPU v5p}（2023 年发布）
单芯片提供 459 TFLOPS 的 BF16 性能，
95GB HBM 内存。
TPU v5p 的核心优势是
\textbf{ICI（Inter-Chip Interconnect）网络}：
一个 TPU v5p pod 可以将 8,960 块芯片通过 ICI 互连，
在 pod 内实现无需显式通信编程的分布式计算。

\textbf{TPU v6e（Trillium）}（2024 年发布）
是效率优化版本，
单芯片性能提升 67\%（相比 v5e），
HBM 容量翻倍。
v6e 特别针对推理场景优化，
是 Gemini 模型的主要推理硬件。

Google 已正式发布 \textbf{TPU v7（代号 Ironwood）}——
这是 TPU 历史上首个以“推理优先”（inference-first）设计理念
打造的版本。
单芯片提供 4,614 TFLOPS 的 FP8 峰值算力
（BF16 为 2,307 TFLOPS；相比 v5p 的 459 TFLOPS 提升了 10 倍），
单个 Ironwood pod 可容纳 9,216 块芯片，
形成一个通过 ICI 高速互连的超大规模推理集群。
Ironwood 的“推理优先”设计意味着
Google 判断推理（而非训练）将成为
未来 AI 计算资源消耗的主要场景——
这与推理时计算（test-time compute）的趋势高度一致。
TPU 的独特优势在于
Google 可以端到端优化
从编译器（XLA）到框架（JAX）到模型架构（Gemini）
到硬件的整个栈，
这种垂直整合在某些场景下
可以提供比通用 GPU 更高的效率。

\subsection{AWS Trainium：云厂商的自研芯片之路}

Amazon 的 Trainium 代表了
云服务商自研 AI 芯片的趋势：

\textbf{Trainium2}（2024 年底推出 Trn2 instances）
提供了相当于 H100 约 30--50\% 的单芯片性能，
但以显著更低的价格提供给 AWS 客户。
Trainium2 的核心价值不在于绝对性能，
而在于\textbf{性价比}——
AWS 可以将自研芯片的成本优势传递给用户。

\textbf{Trainium3}（2025 年 12 月 GA，AWS 的第四代 AI 芯片）
是一个重大升级，
采用 3nm 工艺，是当前\textbf{已量产的最先进制程 AI 芯片}
（与 NVIDIA Blackwell 同代），
单芯片提供 2.52 PFLOPS 的峰值算力，
被 SemiAnalysis 称为
“Amazon Basics 版的 GB200 NVL72”。
Trn3 UltraServers 将 64 块 Trainium3 芯片
通过高速互连组成一个节点，
提供接近 NVIDIA GB200 NVL72 的系统级训练性能，
但价格预计低 30--40\%。
\textbf{Trainium4}（已进入路线图，预计 2027 年）
将进一步提升每芯片的计算密度，
并宣布支持 \textbf{NVLink Fusion}——
允许 Trainium 芯片通过 NVIDIA NVLink 互连与 NVIDIA GPU 混合组网。
这是一个打破传统供应商边界的大胆举措，
意味着未来的 AI 集群可能不再是单一供应商的同构系统。

Anthropic 是 Trainium 的\textbf{锚定客户}（anchor customer），
据报道 Anthropic 已采购约 100 万块 Trainium 芯片。
Claude 模型的部分训练和推理已经运行在 Trainium 上，
这对验证非 NVIDIA 硬件的可行性具有标志性意义。

\begin{keybox}[AI 加速器的竞争格局（2026 年）]
NVIDIA 仍然占据 AI 训练市场约 80\%+ 的份额，
但竞争格局正在发生结构性变化：
（1）AMD 的 MI350 在大内存推理场景中具有竞争力；
（2）Google TPU 通过垂直整合在自家工作负载上效率领先；
（3）AWS Trainium3 以性价比切入中端市场；
（4）国产加速器（华为昇腾 910C、寒武纪等）
在中国市场因出口管制而获得替代需求。
多供应商策略正在成为大型 AI 团队的标准做法，
减少对单一供应商的依赖，
同时通过竞争压低采购成本。
但 NVIDIA 的 CUDA 生态仍然是其最大的护城河，
迁移到其他硬件的软件适配成本是许多团队犹豫的主要原因。
\end{keybox}

\section{国产芯片生态：从替代品到主力军}
\label{sec:domestic-chips}

在全球 AI 加速器竞争加剧的同时，
中国国产 AI 芯片正在经历从“被迫替代”到“主动选择”的转型。
2026 年初的一系列里程碑事件
表明国产芯片生态已经跨越了“可用”的门槛，
正在向“好用”迈进。

\subsection{华为昇腾：中国 AI 芯片的中流砥柱}

\textbf{昇腾 910C}
是当前中国 AI 训练的主力芯片：
提供约 800 TFLOPS 的 FP16 算力，
配备 96GB HBM2e 显存。
虽然在绝对性能上仍与 NVIDIA H100 存在差距
（H100 提供约 989 TFLOPS BF16 dense 和 80GB HBM3），
但 910C 的性能已经足以支撑千亿参数级模型的训练。
据行业估算，2026 年 910C 的计划出货量约为 60 万块，
将成为中国 AI 基础设施的核心供给。

华为的芯片路线图展现了激进的演进节奏：
\textbf{昇腾 910C} → \textbf{昇腾 950PR}
（2026 年 Q1，搭载华为自研 HBM）→
\textbf{昇腾 950DT}（2026 年 Q4）→
\textbf{昇腾 960}（2027 年）→
\textbf{昇腾 970}（2028 年）。
特别值得关注的是 950PR 将首次采用华为自研 HBM，
这意味着华为正在追求从芯片设计到封装存储的
全栈自主可控，减少对 SK 海力士和三星的依赖。

\subsection{“零 NVIDIA” 训练的里程碑}

\textbf{GLM-5}（2026 年 2 月 12 日，智谱 AI）
是中国 AI 发展史上的一个分水岭，
744B 参数的模型完全在 100,000 块华为昇腾 910B 上训练，
\textbf{零 NVIDIA 依赖}。
GLM-5 使用了 28.5 万亿 token 的训练数据，
其性能在多项基准测试上达到了国际一线水平。
这是首次有千亿参数级基础模型完全脱离 NVIDIA 生态完成训练，
证明了国产芯片在大规模训练场景中的工程可行性。

同期发布的 \textbf{GLM-Image}（2026 年 1 月 16 日）
进一步验证了国产芯片在多模态训练中的能力，
它是首个完全在国产芯片上训练的 SOTA 级多模态模型，
使用华为昇腾 Atlas 800T A2 集群和昇思 MindSpore 框架。

\subsection{DeepSeek 的国产芯片战略}

DeepSeek V4 的芯片策略值得特别关注。
据报道，DeepSeek 在 V4 的预发布阶段
将预览访问权限提供给了国产芯片厂商，
但\textbf{刻意排除了 NVIDIA 和 AMD}。
这一决策背后的逻辑是多层次的：
一方面，针对昇腾 910B/C 的专门优化
可以确保模型在国产芯片上的原生性能；
另一方面，这也是对中国 AI 芯片生态的战略支持。

DeepSeek 还提供了便捷的国产芯片适配方案：
通过其开源工具链，开发者只需\textbf{一行代码}
即可将 CUDA 调用转换为华为 CANN 调用——
大幅降低了从 NVIDIA GPU 迁移到昇腾的门槛。

\subsection{其他国产芯片力量}

\textbf{寒武纪}（Cambricon）
是另一支重要的国产 AI 芯片力量，
2026 年的芯片出货目标为 50 万块，
公司估值已达约 880 亿美元（约 6,400 亿人民币）。
寒武纪在推理加速器领域尤其有竞争力，
多家中国互联网公司已采用寒武纪芯片进行推理部署。

\begin{warnbox}[国产芯片的核心挑战：长周期训练稳定性]
尽管国产芯片在单芯片性能和集群规模上已取得显著进步，
但\textbf{长周期训练稳定性}仍然是最大的技术挑战。
训练一个千亿参数模型通常需要数月时间，
在此期间任何硬件故障、驱动 bug 或数值不稳定
都可能导致训练中断和 checkpoint 回滚。
NVIDIA 在这方面有十余年的工程积累和广泛的用户验证，
而国产芯片生态的成熟度仍在快速追赶中。
GLM-5 在十万块 910B 上成功完成训练，
正是对这一挑战的有力回答，
但将这种能力从少数头部团队推广到更广泛的用户群体，
仍需要整个生态系统的持续投入。
\end{warnbox}

\section{端侧模型：让 AI 无处不在}

不是所有场景都需要云端的 700B 模型。
端侧（on-device）AI 正在成为一个重要的发展方向。

\subsection{为什么端侧模型重要？}

端侧推理的核心价值：

\begin{itemize}[noitemsep]
  \item \textbf{隐私}：数据不离开设备，
        对于医疗记录、私人通信、企业机密等敏感场景至关重要。
  \item \textbf{延迟}：本地推理消除了网络往返时间，
        首 token 延迟可以低至数十毫秒。
  \item \textbf{离线能力}：飞行模式、地铁、偏远地区，
        不依赖网络连接的 AI 能力。
  \item \textbf{成本}：没有 API 调用费用。
        对于高频率使用场景（如写作助手、代码补全），
        本地模型的边际成本几乎为零。
\end{itemize}

\subsection{NPU 硬件的快速演进}
\label{sec:npu-hardware}

端侧 AI 的能力边界由硬件决定。
2025--2026 年，主要芯片厂商的 AI 加速能力正在快速提升：

\textbf{Apple Neural Engine}（ANE）：
Apple 的 M 系列和 A 系列芯片
内置的 Neural Engine 是端侧 AI 的标杆。
M4 芯片的 ANE 提供 38 TOPS 的算力，
配合 Apple Silicon 的统一内存架构
（CPU/GPU/NPU 共享内存，避免数据搬运），
可以在 iPhone 和 Mac 上流畅运行 3B 模型。
Apple Intelligence 基于这一硬件能力构建，
在设备端运行一个 $\sim$3B 的基础语言模型~\cite{gunter2024appleintelligence}。

\textbf{Qualcomm Hexagon NPU}（Snapdragon X Elite/X2 系列）：
Snapdragon X Elite 的 Hexagon NPU 提供 45 TOPS，
是 Windows on ARM PC 上的主要 AI 加速器。
Snapdragon X2（预计 2026 年发布）
进一步提升 NPU 性能，
并改善了对主流 AI 框架的兼容性。
Qualcomm 的 AI Stack（QNN SDK）
支持将 ONNX 和 PyTorch 模型
直接部署到 Hexagon NPU 上。

\textbf{Intel Lunar Lake / Arrow Lake}：
Intel 的最新客户端 CPU 集成了
提供 48 TOPS 的 NPU，
但其软件生态（OpenVINO）
相比 Apple 和 Qualcomm 仍有差距。

\textbf{AMD Ryzen AI}：
AMD 的 Ryzen AI 300 系列集成了
XDNA 2 NPU，提供 50 TOPS，
在 Windows AI PC 市场与 Qualcomm 和 Intel 竞争。

\begin{table}[H]
\centering\small
\renewcommand{\arraystretch}{1.3}
\caption{主要端侧 AI 加速器对比（2026 年初）}
\begin{tabularx}{\linewidth}{l l l >{\raggedright\arraybackslash}X}
\toprule
\rowcolor{figLightGray}
\textbf{芯片} & \textbf{NPU 算力} & \textbf{内存} & \textbf{可运行的模型规模} \\
\midrule
Apple M4 (ANE) & 38 TOPS & 16--32GB 统一 & 3B 流畅，7B 可运行 \\
Snapdragon X Elite & 45 TOPS & 16--32GB LPDDR5x & 3B 流畅，7B 可运行 \\
Intel Lunar Lake & 48 TOPS & 16--32GB LPDDR5x & 3B 流畅 \\
AMD Ryzen AI 300 & 50 TOPS & 16--64GB DDR5 & 3B 流畅，7B 可运行 \\
Apple M4 Max (ANE) & 38 TOPS & 36--128GB 统一 & 7B--14B 流畅 \\
\bottomrule
\end{tabularx}
\end{table}

\subsection{小模型家族与端侧能力}

\textbf{小模型家族}正在快速发展：
Apple Intelligence~\cite{gunter2024appleintelligence}
将一个 $\sim$3B 的 LLM 部署到 iPhone 上；
Google Gemini Nano（1.8B 和 3.25B）专为 Pixel 手机设计；
LLaMA 3.2~\cite{dubey2024llama3} 提供 1B 和 3B 的端侧版本；
Qwen3-0.6B~\cite{qwen2025qwen3} 和 Gemma 3-1B
更是将模型压缩到了 1B 以下。

量化技术（如 GGUF Q4\_K\_M）和
推理优化（如 llama.cpp 的 ARM NEON 和 Metal 优化）
使得在消费级硬件上运行 7B--14B 的模型成为现实。

\textbf{2026 年春季的端侧模型新进展}进一步扩展了这一趋势：
\textbf{Qwen 3.5 小模型系列}（2026 年 3 月 3 日）
开源了 0.8B/2B/4B/9B 四个规格，
均以 Apache 2.0 许可证发布，
特别针对端侧推理场景优化。
\textbf{阶跃星辰 Step 3.5 Flash}（2026 年 2 月 2 日）
的 196B/11B MoE 架构展示了另一种可能，
通过极致的 MoE 稀疏化（仅 5.6\% 的参数活跃），
将原本需要数据中心级硬件的推理能力
压缩到单块 RTX 4090（24GB）上运行。
更值得关注的是，
arXiv:2512.15943 的研究表明，
在 Agent 工具调用（tool calling）场景中，
精心微调的小型 LLM 可以\textbf{超越}大型模型的表现，
这意味着端侧模型在 Agent 生态中
可能扮演比传统认知更重要的角色。

端侧模型的核心能力保留策略包括：
\textbf{知识蒸馏}——用大模型的输出训练小模型，
使小模型「继承」大模型的推理能力
（DeepSeek-R1-Distill 系列就是这种方法的典范）；
\textbf{结构化剪枝}——去除模型中对质量贡献最小的神经元或注意力头，
在保持架构兼容性的同时减小模型；
\textbf{共享词嵌入}——让输入嵌入层和输出层共享权重，
减少参数量（LLaMA 3 等模型采用了这种设计）。

实际测试表明，当前的 3B 模型在以下任务上已经「够用」：
简单问答、文本摘要（短文）、情感分析、
文本分类、短文翻译、格式转换。
它们在以下任务上仍然不足：
复杂推理、长文理解、代码生成（复杂逻辑）、
创意写作（高质量）、多轮对话（需要长期记忆）。

\subsection{端云协作：混合推理}

端侧--云端的\textbf{混合推理}（hybrid inference）
可能是最具实用价值的方向：
简单请求在设备端本地处理（零延迟、完全隐私），
复杂请求上传到云端的大模型处理（更高质量）。

Apple Intelligence 的架构正是这种模式的典范：
端侧 3B 模型处理大部分请求，
当模型判断自己无法给出高质量回答时，
请求被路由到云端更大的模型
（Apple 称之为 Private Cloud Compute，
使用专用的 Apple Silicon 服务器，
承诺数据不存储、不记录）。
路由决策可以由端侧模型自己做出，
它可以评估自己对回答的「信心」，
信心不足时自动升级到云端。

这种架构的经济优势明显：
假设 80\% 的请求可以在端侧处理，
那么云端的服务器需求减少 80\%，
显著降低了 API 提供商的运营成本。

\begin{corebox}[2026 年端侧 AI 的能力阈值]
NPU 硬件的发展正在快速改变端侧 AI 的能力阈值。
2024 年的 iPhone 可以流畅运行 3B 模型，
2026 年的旗舰手机（搭载 A19 或 Snapdragon X2）
预计可以运行 7B 模型。
高端笔记本（搭载 M4 Max/Ultra，128GB 统一内存）
已经可以流畅运行 70B 的量化模型。
这意味着一些原本只能在服务器上运行的任务：
如本地代码补全、文档摘要、个性化写作助手：
将逐步转移到用户设备上，
从根本上改变 AI 服务的交付模式。
\end{corebox}

\section{AI 监管：规则正在形成}
\label{sec:regulation}

随着 LLM 能力的快速提升，
全球各主要经济体正在加速构建 AI 监管框架。
监管环境对 AI 公司的战略选择、
开源生态和技术发展方向有深远影响。

\subsection{EU AI Act：全球最严格的 AI 法规}

欧盟 AI 法案（EU AI Act）于 2024 年 8 月 1 日正式生效，
是全球首个全面的 AI 监管法律框架。
其核心是\textbf{风险分级}制度：

\begin{itemize}[noitemsep]
  \item \textbf{不可接受风险}（禁止）：
        社会评分系统、实时远程生物识别（执法例外）、
        操纵性 AI 系统。
  \item \textbf{高风险}（严格监管）：
        用于关键基础设施、教育、就业、司法的 AI 系统。
        需要进行一致性评估（conformity assessment）、
        建立风险管理系统、确保数据治理、保留技术文档。
  \item \textbf{通用目的 AI 模型}（GPAI，特别规定）：
        所有 GPAI 提供商必须提供技术文档、
        遵守版权法、发布训练数据摘要。
        “systemic risk” 模型（训练算力超过 $10^{25}$ FLOPs）
        还需进行对抗性测试、监控严重事件、
        确保网络安全。
  \item \textbf{有限风险 / 最小风险}（轻度或无监管）：
        聊天机器人、AI 生成内容等，
        仅需透明度标注（告知用户正在与 AI 交互）。
\end{itemize}

\textbf{关键时间节点}：
2025 年 2 月，禁止不可接受风险 AI 的条款生效；
2025 年 8 月，GPAI 模型条款和治理结构生效；
2026 年 8 月，高风险 AI 系统的全面合规要求生效。
违规罚款最高可达 3,500 万欧元
或全球年营业额的 7\%（取较高者）。

EU AI Act 对开源模型有一定的豁免，
“free and open-source” 的 GPAI 模型
如果不构成 systemic risk，
可以免除部分合规义务（如技术文档要求的简化），
但仍需遵守版权法。

\subsection{美国：碎片化但加速中的监管}

美国没有联邦级别的统一 AI 法律，
但多层次的监管正在形成：

\textbf{联邦层面}：
拜登政府 2023 年 10 月的行政令
（Executive Order on AI Safety and Security）
要求开发 dual-use 基础模型的公司
向联邦政府报告训练细节和安全测试结果。
特朗普政府在 2025 年 1 月上任后撤销了这一行政令，
转向更为宽松的 “pro-innovation” 政策。
但国会层面，
多项 AI 立法提案正在推进中，
包括 AI 系统的透明度要求、
深度伪造标注义务和 AI 在关键领域的使用规范。

\textbf{州层面}：
加利福尼亚州的 SB 1047（AI Safety Bill）
在 2024 年引发了激烈辩论，
该法案要求 AI 开发者对模型造成的严重危害承担责任，
最终被州长 Gavin Newsom 否决，
理由是过于严厉可能遏制创新。
但 2025--2026 年的后续立法尝试仍在继续。
科罗拉多州于 2024 年通过了
美国首个全面的 AI 算法歧视法律。

\textbf{机构层面}：
NIST 的 AI Risk Management Framework（AI RMF）
提供了自愿性的 AI 风险管理指南；
FTC 将 AI 虚假声明纳入消费者保护执法范围；
SEC 开始关注 AI 在金融市场中的风险。

\subsection{中国：审批制 + 内容管控}

中国的 AI 监管采取了\textbf{分类审批}的模式：

\begin{itemize}[noitemsep]
  \item \textbf{生成式 AI 管理办法}（2023 年 8 月生效）：
        向公众提供生成式 AI 服务
        需要通过算法备案和安全评估。
        服务提供者必须确保生成内容
        符合社会主义核心价值观，
        不得生成「违法」内容。
  \item \textbf{深度合成管理规定}（2023 年 1 月生效）：
        AI 生成的图像、音频、视频必须标注。
  \item \textbf{算法推荐管理规定}（2022 年 3 月生效）：
        覆盖推荐算法的透明度和用户控制权。
\end{itemize}

中国的监管特点是\textbf{事前审批}（pre-deployment approval）——
模型必须通过政府审查才能上线，
这与欧美的事后监管（post-deployment regulation）形成对比。
截至 2025 年底，已有超过 200 个 AI 大模型
通过了中国的备案审批。
这种制度对开源模型的影响是复杂的，
开源模型的权重可以自由下载，
但基于开源模型构建的面向公众的服务
仍然需要通过审批。

\begin{warnbox}[监管对 AI 发展的影响]
AI 监管正在形成一个全球「三极格局」：
\textbf{欧盟}——严格的规则驱动，侧重权利保护和透明度，
可能减慢创新速度但提高公众信任；
\textbf{美国}——碎片化但灵活，侧重创新友好，
可能导致监管套利和执法不一致；
\textbf{中国}——集中审批制，侧重内容管控和国家安全，
在保证可控性的同时可能限制学术自由度。
对于 AI 公司来说，合规成本正在成为
不可忽视的运营开支，
特别是需要在多个司法管辖区运营的全球化公司，
需要同时满足不同且有时矛盾的监管要求。
开源模型在这个格局中处于特殊位置：
一方面，开源天然满足透明度要求；
另一方面，开源后的模型无法被收回，
一旦被用于违规目的，责任归属变得模糊。
\end{warnbox}

\section{AI 与社会：机遇、冲击与应对}
\label{sec:ai-and-society}

LLM 能力的快速提升不仅是技术事件，
更是深刻的社会事件。
2025--2026 年，AI 对就业、教育、创意产业
和社会结构的影响从理论预测变成了可观察的现实。

\subsection{就业市场的结构性变化}

\textbf{Goldman Sachs 的估计}（2023 年 3 月发布的报告，
2025 年更新）预测 AI 可能影响全球约 3 亿个全职工作岗位，
其中约 25\% 的工作任务可以被当前的生成式 AI 完成。
但“影响”不等于“替代”，
报告的核心观点是大多数职业将经历\textbf{任务重组}
（task restructuring）而非完全消失：
一些任务被自动化，
释放出来的时间被用于 AI 无法胜任的高价值活动。

到 2026 年初，实际数据开始验证和修正这些预测：

\textbf{软件工程}是受影响最直接的行业。
多项调查显示，使用 AI 编程工具的开发者
报告生产力提升 30--55\%（GitHub 的内部研究为 55\%，
Google 的内部研究为 30\%）。
然而，“生产力提升”并未导致大规模裁员，
相反，软件工程师的需求在 2025 年保持稳定甚至增长，
因为 AI 辅助使得更多项目变得经济可行，
\textbf{扩大了软件开发的总市场规模}。
但入门级编程岗位的需求确实出现了下降，
AI 正在承担越来越多的初级编码工作，
使得企业对低经验开发者的需求减少。

\textbf{内容创作}经历了最痛苦的调整。
翻译行业在 2024--2025 年遭受了严重冲击，
多家大型翻译公司报告业务量下降 30--50\%，
特别是标准化的技术文档翻译几乎完全被 AI 取代。
但高端翻译（文学翻译、法律文件、外交文本）
因需要深层文化理解和责任承担而保持需求。
类似的模式在插画、平面设计、广告文案等领域重现，
标准化的、模板驱动的工作被 AI 快速侵蚀，
但需要独创性和品牌理解的高端工作
反而因为 AI 降低了入门门槛而面临更多竞争。

\textbf{知识工作}（法律、咨询、金融分析）的变化更为微妙。
法律领域的文档审查和合同分析
已经被 AI 显著加速（报告效率提升 3--5 倍），
但律师岗位的总数并未明显减少，
因为加速的文档处理释放出的时间
被重新分配到更高价值的策略咨询和客户关系管理上。
McKinsey 的 2025 年报告估计，
法律和咨询行业将经历约 15\% 的任务自动化，
但仅导致约 5\% 的岗位净减少。

\subsection{“暂停”辩论与安全组织}

2023 年 3 月，Future of Life Institute 发布
\textbf{“暂停巨型 AI 实验”} 公开信，
征集到超过 30,000 个签名
（包括 Elon Musk、Steve Wozniak、
Yoshua Bengio 等知名人士），
呼吁所有 AI 实验室暂停训练
比 GPT-4 更强大的 AI 系统至少 6 个月。

这封信引发了 AI 安全辩论的空前激烈化。
支持者认为，AI 能力的发展速度
已经超越了人类对其风险的理解速度，
暂停是负责任的预防措施。
反对者，包括 Yann LeCun 和多数开源社区——
认为暂停在实践中不可执行
（无法验证各国和企业是否遵守），
且可能将技术领先地位拱手让给不遵守暂停的竞争对手。

事实上，暂停从未发生。
但这场辩论催化了多个重要的安全组织和倡议：

\textbf{Anthropic} 自称为“AI safety company”，
其 Responsible Scaling Policy（RSP）
定义了一系列能力阈值（ASL-1 到 ASL-4），
当模型能力达到特定阈值时
触发相应的安全评估和部署限制。
ASL-3（“能够显著辅助大规模攻击的能力”）
是当前最紧迫的关注点，
Claude 4 系列的安全评估正是基于 ASL-3 标准进行的。

\textbf{OpenAI} 成立了安全顾问委员会（Safety Advisory Committee），
发布了 Preparedness Framework，
定义了模型在网络安全、CBRN（化学、生物、放射、核）、
说服操纵和模型自主性四个维度上的风险等级。
然而，2024 年 Ilya Sutskever 和 Jan Leike
从 OpenAI 的 Superalignment 团队离职
引发了外界对 OpenAI 安全承诺的质疑，
批评者认为 OpenAI 在追求能力提升和安全保障之间
存在结构性的利益冲突。

\textbf{AI Safety Institute}（英国，2023 年成立；
美国，2024 年成立）
是政府层面的安全研究机构，
负责对 frontier 模型进行独立的安全评估。
英国 AISI 已经与 OpenAI、Anthropic、Google 和 Meta
签署了模型发布前的安全测试协议。

\subsection{教育系统的重构}

AI 对教育系统的冲击是多维度的。
学生使用 ChatGPT 写作业已经从“作弊丑闻”
演变为“教育系统需要应对的新现实”。
2025 年，多数主要大学已经从“禁止使用”
转向“教学生如何负责任地使用 AI 工具”的策略。

更深层的影响在于\textbf{教育目标的重新定义}。
当 AI 可以生成流利的论文、解答数学题、编写代码时，
“掌握知识”和“执行标准化操作”的价值下降了。
取而代之的是\textbf{批判性思维}
（判断 AI 输出的正确性和适当性）、
\textbf{提问能力}
（知道问什么比知道答案更重要）
和\textbf{跨学科整合}
（AI 擅长单一领域的深入分析，
但将不同领域的知识创造性地连接仍然是人类的优势）。

\begin{warnbox}[AI 与社会的“不对称冲击”]
AI 对社会的冲击是高度不对称的。
\textbf{技能溢价加剧}：
能够有效利用 AI 工具的知识工作者
看到了生产力和收入的显著提升，
而无法适应的工作者面临技能贬值。
\textbf{地理不对称}——
AI 产业高度集中在美国、中国和少数欧洲国家，
发展中国家面临被进一步边缘化的风险。
\textbf{代际不对称}——
年轻一代更快地适应 AI 工具，
但也更可能面临入门级岗位被自动化的困境。
应对这种不对称冲击需要的不仅是技术进步，
更是教育改革、社会安全网重构
和国际技术合作框架的建立。
Anthropic CEO Dario Amodei 在
“Machines of Loving Grace”（2025 年）中
勾勒了 AI 增强人类而非替代人类的愿景，
但实现这一愿景需要有意识的制度设计，
而非仅靠市场的自发调节。
\end{warnbox}

%% --- TikZ: Frontier Capability Radar Chart ---
\begin{figure}[htbp]
\centering
\begin{tikzpicture}[
  every node/.style={font=\small},
  >={Stealth[length=2.5pt,width=2.5pt]},
]

% --- Constants ---
\def\R{3.2}  % max radius

% --- Concentric octagonal grid ---
\foreach \frac/\pct in {0.25/25, 0.50/50, 0.75/75, 1.0/100} {
  \draw[figLightGray, line width=0.4pt]
    (90:{\frac*\R})  -- (45:{\frac*\R})  -- (0:{\frac*\R})
    -- (-45:{\frac*\R})  -- (-90:{\frac*\R})  -- (-135:{\frac*\R})
    -- (180:{\frac*\R})  -- (135:{\frac*\R})  -- cycle;
}
% Percentage labels along 90-degree axis
\foreach \frac/\pct in {0.25/25, 0.50/50, 0.75/75, 1.0/100} {
  \node[font=\tiny, figGray!60, anchor=south west] at (90:{\frac*\R}) {\pct\%};
}

% --- Radial axes ---
\foreach \ang in {90, 45, 0, -45, -90, -135, 180, 135} {
  \draw[figGray!30, line width=0.3pt] (0,0) -- (\ang:\R);
}

% --- 2024 polygon (behind, dashed, gray) ---
\fill[figGray!8, fill opacity=0.6]
  (90:{0.50*\R}) -- (45:{0.38*\R}) -- (0:{0.25*\R})
  -- (-45:{0.63*\R}) -- (-90:{0.25*\R}) -- (-135:{0.13*\R})
  -- (180:{0.38*\R}) -- (135:{0.50*\R}) -- cycle;
\draw[figGray!60, line width=0.6pt, dashed]
  (90:{0.50*\R}) -- (45:{0.38*\R}) -- (0:{0.25*\R})
  -- (-45:{0.63*\R}) -- (-90:{0.25*\R}) -- (-135:{0.13*\R})
  -- (180:{0.38*\R}) -- (135:{0.50*\R}) -- cycle;

% --- 2026 polygon (front, solid, blue) ---
\fill[figBlue!10, fill opacity=0.5]
  (90:{0.95*\R}) -- (45:{0.88*\R}) -- (0:{0.80*\R})
  -- (-45:{0.95*\R}) -- (-90:{0.63*\R}) -- (-135:{0.38*\R})
  -- (180:{0.70*\R}) -- (135:{0.88*\R}) -- cycle;
\draw[figBlue!70, line width=0.6pt]
  (90:{0.95*\R}) -- (45:{0.88*\R}) -- (0:{0.80*\R})
  -- (-45:{0.95*\R}) -- (-90:{0.63*\R}) -- (-135:{0.38*\R})
  -- (180:{0.70*\R}) -- (135:{0.88*\R}) -- cycle;

% --- Vertex dots ---
\foreach \ang/\frac in {90/0.95, 45/0.88, 0/0.80, -45/0.95,
                         -90/0.63, -135/0.38, 180/0.70, 135/0.88} {
  \fill[figBlue!70] (\ang:{\frac*\R}) circle (1.2pt);
}
\foreach \ang/\frac in {90/0.50, 45/0.38, 0/0.25, -45/0.63,
                         -90/0.25, -135/0.13, 180/0.38, 135/0.50} {
  \fill[figGray!50] (\ang:{\frac*\R}) circle (1pt);
}

% --- Axis labels (offset to avoid overlap) ---
\node[font=\scriptsize\bfseries, figGray, above=4pt]       at (90:\R)   {推理};
\node[font=\scriptsize\bfseries, figGray, above right=3pt]  at (45:\R)   {多模态};
\node[font=\scriptsize\bfseries, figGray, right=4pt]        at (0:\R)    {Agent};
\node[font=\scriptsize\bfseries, figGray, below right=3pt]  at (-45:\R)  {代码};
\node[font=\scriptsize\bfseries, figGray, below=4pt]        at (-90:\R)  {科学};
\node[font=\scriptsize\bfseries, figGray, below left=3pt]   at (-135:\R) {具身AI};
\node[font=\scriptsize\bfseries, figGray, left=4pt]         at (180:\R)  {端侧};
\node[font=\scriptsize\bfseries, figGray, above left=3pt]   at (135:\R)  {长上下文};

% --- Legend ---
\draw[figGray!60, line width=0.6pt, dashed] (-1.2, -4.2) -- (-0.5, -4.2);
\fill[figGray!8] (-1.2, -4.35) rectangle (-0.5, -4.05);
\node[font=\scriptsize, figGray, right=2pt] at (-0.5, -4.2) {2024 年};

\draw[figBlue!70, line width=0.6pt] (1.0, -4.2) -- (1.7, -4.2);
\fill[figBlue!10] (1.0, -4.35) rectangle (1.7, -4.05);
\node[font=\scriptsize, figBlue, right=2pt] at (1.7, -4.2) {2026 年};

\end{tikzpicture}
\caption{Frontier 模型能力雷达图：2024 年 vs.\ 2026 年。
推理、多模态、代码能力在两年内经历了从“中级”到“顶级”的飞跃，
而具身 AI 和 AI for Science 仍处于早期发展阶段。}
\label{fig:frontier-radar}
\end{figure}


\section{开合之争：开源 vs.\ 闭源的哲学与实践}

2025--2026 年，开源与闭源模型之间的「开合之争」
是 AI 领域最重要的叙事之一。

\subsection{闭源的论据}

\begin{itemize}[noitemsep]
  \item \textbf{安全控制}：
        模型权重不公开，可以通过 API 层面的安全措施
        （内容过滤、速率限制、使用监控）
        防止滥用。一旦权重开源，安全对齐可以被移除。
  \item \textbf{竞争优势}：
        训练一个 frontier model 需要数千万到数亿美元的投入，
        开源意味着竞争对手可以免费使用你的成果。
  \item \textbf{责任管理}：
        如果模型被用于有害目的，
        API 提供商可以追踪和阻止滥用者；
        开源后则无法控制模型的使用方式。
  \item \textbf{收入模型}：
        API 定价（按 token 收费）
        是当前 LLM 公司最直接的变现方式。
\end{itemize}

\subsection{开源的论据}

\begin{itemize}[noitemsep]
  \item \textbf{创新速度}：
        开源社区的创新速度远超任何单一公司。
        LoRA~\cite{hu2021lora}、QLoRA~\cite{dettmers2023qlora}、
        Flash Attention~\cite{dao2022flashattention} 等
        关键创新都来自开源社区或学术界。
  \item \textbf{科学可重复性}：
        如果研究成果无法被其他团队验证和复现，
        其科学价值就打了折扣。
  \item \textbf{民主化}：
        防止 AI 能力集中在少数公司手中。
        学术研究者、初创企业、发展中国家的开发者
        都应该能够获取和使用先进的 AI 技术。
  \item \textbf{国家安全}：
        一个国家的 AI 能力不应该依赖于外国公司的 API，
        开源模型允许本地部署和主权控制。
  \item \textbf{生态系统效应}：
        当你的模型被广泛使用，
        围绕它建立的工具、教程、微调模型和应用
        构成了一个自我强化的生态系统。
\end{itemize}

\subsection{“开放权重” 不等于真正的开源}

一个重要的区分：
当前所谓的「开源模型」大多只是 \textbf{open weights}——
公开了模型权重，但训练数据、完整的训练代码、
数据处理管道和评估基础设施通常不公开。

从技术角度看，
开源模型的质量已经接近甚至达到闭源模型的水平，
DeepSeek-R1 在数学推理上与 o1 相当，
LLaMA 4 Maverick 在综合能力上逼近 GPT-4o。
这是因为 LLM 的核心技术
（Transformer 架构、训练方法、数据处理流水线）
已经高度标准化，
真正的差异化越来越依赖于\textbf{工程执行}而非秘密技术。

一个更精确的开放程度光谱是：
\textbf{完全闭源}（仅内部使用）→
\textbf{API-only}（通过 API 提供服务，不公开权重）→
\textbf{开放权重}（公开权重，不公开训练代码/数据）→
\textbf{开放训练代码}（公开权重和训练代码，不公开数据）→
\textbf{完全开源}（数据 + 代码 + 权重全部公开）。

目前只有极少数模型（如 OLMo by AI2、SmolLM2 by HuggingFace~\cite{allal2025smollm2}）
处于「完全开源」的端点。
DeepSeek-V3/R1 虽然使用 MIT 许可证，
但其训练数据和完整的训练基础设施并未公开。

这种区分不仅是语义学上的挑剔，
它有\textbf{实际的技术后果}。
仅有权重的“开源”模型
可以被微调和部署，但无法被\textbf{复现}：
如果不知道训练数据和训练流程的细节，
其他团队无法验证报告的性能数字，
也无法理解模型的弱点和偏见来源。
这在安全评估中尤其重要：
如果训练数据中包含有毒内容，
仅通过黑盒测试很难全面发现所有问题。

\textbf{OSI（Open Source Initiative）}在 2024 年发布了
\textbf{Open Source AI Definition}（OSAID 1.0），
明确定义了“开源 AI”的四项要求：
（1）使用模型进行任何目的，无需请求许可；
（2）研究模型的工作方式并检查其组件；
（3）修改模型以适应任何目的；
（4）分享模型供他人使用。
根据这一定义，
大多数“开放权重”模型因缺少训练数据和完整代码
而\textbf{不符合}“开源 AI”的标准，
它们更准确地应该被称为“available-weight models”。
这一定义的发布在社区中引发了激烈辩论，
但它至少为这场讨论提供了一个更精确的参照框架。

\subsection{Meta 的开源战略}

Meta~\cite{touvron2023llama,dubey2024llama3} 的开源决策
是 AI 行业最值得研究的战略案例之一。
Mark Zuckerberg 的逻辑是\textbf{商品化互补品}（commoditize the complement）：
Meta 的核心业务是社交网络和广告，
AI 模型是这些业务的基础设施。
如果 AI 模型成为商品（人人都能使用），
那么构建在 AI 之上的应用
（如 Meta 的推荐系统、广告排序、内容审核）
才是真正的差异化竞争力。
开源 LLaMA 还能吸引顶尖人才，
研究者更愿意为一个开放的项目工作，
因为他们的工作能被全世界看到。

\subsection{DeepSeek-R1 的冲击}

DeepSeek-R1~\cite{deepseekr1} 的发布
深刻改变了开源 vs.\ 闭源的辩论格局。
一个来自中国的、团队规模不到 200 人的实验室，
训练出了在数学和代码推理上与 GPT-4 级别模型匹敌的模型，
并以 MIT 许可证完全开源。

这对闭源阵营的打击是多维度的：
（1）如果开源模型能达到这个水平，
那么闭源的「技术壁垒」论据就被削弱了；
（2）DeepSeek-V3 报告的最终预训练 GPU 租赁成本约 \$5.5M
（不含前期研发、数据处理和后训练成本），
远低于 GPT-4 传闻中的 \$100M+（该数据未经官方确认），
这说明 frontier model 的训练不一定需要巨额投入；
（3）开源模型的性能达到商用水平后，
API 定价将面临巨大的下行压力。

\subsection{2026 年春季：开源的“定义性时刻”}

如果说 DeepSeek-R1 是开源的“第一枪”，
那么 2026 年春季的模型发布潮则是“全面突破”。
几乎所有中国头部实验室在这一波发布中
都选择了\textbf{开放权重}发布：
GLM-5（744B，MIT 许可证）、
Kimi K2.5（1T MoE，Apache 2.0）、
Qwen 3.5（397B/17B，Apache 2.0）、
阶跃星辰 Step 3.5 Flash（196B/11B，MIT 许可证）。

这带来了几个深远的影响：

\textbf{中国开源的全球影响力}：
中国模型在 HuggingFace 上的使用占比
从 2024 年的 1.2\% 飙升至 2026 年初的 \textbf{30\%}——
这是一个数量级的跳跃。
中文 NLP 不再是开源社区的边缘话题，
而是与英文模型并驾齐驱的主流力量。

\textbf{价格战加速智能民主化}：
开源模型的性能追平闭源模型后，
API 定价发生了剧烈的下行压力。
MiniMax 的 API 定价仅为 GPT-5 的 1/10 到 1/20；
豆包的定价比 GPT-5.2 便宜 3.7--5.9 倍；
Qwen 的 API 比同级别闭源模型便宜 60\%。
“Intelligence too cheap to meter”
（智能便宜到不值得计量）
这句借自核能时代的预言正在成为现实，
对于大多数应用场景，LLM 调用的成本
已经不再是阻碍采用的因素。

\section{Scaling 的尽头与新方向}
\label{sec:scaling-new-directions}

“Scaling Laws”（缩放定律）~\cite{kaplan2020scaling,hoffmann2022chinchilla}
一直是 LLM 进步的核心驱动力，
更多的数据、更大的模型、更多的计算，
几乎确定性地带来更好的性能。
但 2025--2026 年的多个信号表明，
\textbf{传统的 pre-training scaling 可能正在接近瓶颈}——
新的 scaling 方向正在涌现。

\subsection{预训练 Scaling 的放缓信号}

\textbf{数据墙}是最直接的瓶颈。
人类历史上积累的高质量文本数据
，科学论文、书籍、高质量网页、代码仓库，
是有限的。据 Epoch AI 的估计~\cite{villalobos2024position}，
互联网上的高质量英文文本总量约为 4--5 万亿 token。
考虑到中文和其他语言的数据，
总量可能达到 10--15 万亿 token。
而 GPT-5 和 Gemini 3 的训练数据量已经接近这个上限。

应对数据墙的策略包括：
\textbf{合成数据}——用大模型生成训练小模型或自身下一版本的数据
（但存在“模型坍缩”风险，
模型在自身生成的数据上训练可能导致分布退化~\cite{shumailov2024curse}）；
\textbf{数据混合优化}——
同样的数据，不同的混合比例和课程顺序
可以带来显著不同的训练效果
（DeepSeek 在这方面有深入的工程积累）；
\textbf{多模态数据}——
图像、视频和音频数据的总量远超纯文本数据，
这也是原生多模态成为主流的驱动力之一。

\textbf{算力-性能边际递减}。
Scaling Laws 预测 loss 随计算量
以\textbf{幂律}（power law）形式下降，
$L \propto C^{-\alpha}$，
其中 $\alpha \approx 0.05$（Chinchilla 估计~\cite{hoffmann2022chinchilla}）。
这意味着将 loss 降低 10\%
需要将计算量增加约 $10^{1/\alpha} \approx$ 一个数量级。
当训练成本已经达到数十亿美元时，
下一个 10\% 的改善需要数百亿美元的投入，
经济上的可持续性成为核心约束。

\subsection{新的 Scaling 维度}

面对预训练 scaling 的瓶颈，
AI 领域正在探索\textbf{多维度的 scaling 策略}：

\textbf{推理时 Scaling（Test-time Compute Scaling）}。
这是 2025--2026 年最重要的新 scaling 维度——
不增大模型，而是在推理时投入更多计算。
o1/o3 系列和 DeepSeek-R1 证明了
模型可以通过更长的推理链、
更多的候选方案探索和自我验证
在固定参数下显著提升输出质量。
o4-mini 的成功尤其说明了这一点：
一个中等规模的模型通过充分的推理时计算
可以匹配甚至超越更大的直接生成模型。

\textbf{架构效率 Scaling}。
MoE 架构代表了另一种 scaling 哲学——
不是让所有参数都参与每次推理，
而是通过路由机制只激活一小部分参数。
DeepSeek-V3 的 671B 总参数 / 37B 活跃参数
意味着仅约 5.5\% 的参数被激活。
Step 3.5 Flash 将这一比例推向了极致的 5.6\%。
这等效于\textbf{在固定的推理成本下
扩大了模型可以“查阅”的知识库}。

\textbf{记忆 Scaling}。
DeepSeek 的 Engram 论文提出了第三种维度，
不增加计算参数，而是扩展模型的\textbf{记忆容量}。
Engram 的 U 型缩放定律表明，
在固定总参数预算下，
将 75\% 的参数分配给记忆（$O(1)$ 查找）、
仅 25\% 分配给计算（$O(n^2)$ 注意力）
可以达到最优性能。
如果这一发现在更大规模上得到验证，
它可能从根本上改变模型架构的设计哲学：
从“更多的注意力层”转向“更大的知识存储 + 精简的推理引擎”。

\textbf{工具 Scaling}。
Agent 范式提供了又一种 scaling 方式，
不增大模型本身的能力，
而是\textbf{扩展模型可以调用的工具}。
一个能够使用计算器的 3B 模型
在数学运算上可以超越一个 70B 的纯文本模型。
MCP 生态系统的爆发式增长
使得“工具 scaling”从理论变成了实践，
截至 2026 年 3 月，超过 10,000 个公共 MCP 服务器
为 Agent 提供了从数据库查询到 API 调用
到文件操作的广泛能力。

\begin{whybox}[Scaling 的终极问题]
一个深层的哲学问题是：
\textbf{预测下一个 token 的目标函数本身是否有上限？}
Scaling Laws 告诉我们，
在给定目标函数下，
性能随计算量的增加而持续改善。
但下一个 token 预测的最优解
是训练数据分布的精确拟合，
一旦模型完美地学会了预测
人类文本中的下一个 token，
它的性能就达到了由训练数据质量决定的天花板。
超越这个天花板需要的不是更多的计算，
而是\textbf{更好的数据}或\textbf{更好的目标函数}。
DeepSeek 的 GRPO 和 Engram
可以被视为在目标函数维度上的探索，
前者让模型学会“如何思考”，
后者让模型学会“如何记忆”，
这些能力不是纯粹的下一个 token 预测能诱导出来的。
\end{whybox}


\section{总结与展望}

2025--2026 年的 LLM 领域
正在经历多条技术和产业路线的同时演进，
且竞争烈度达到了前所未有的水平：

\textbf{推理能力的系统化}：
从提示工程到专用推理模型再到可配置推理预算，
模型正在学会「思考」而非仅仅「回忆」。
test-time compute 正在成为
与训练 compute 同等重要的资源维度。
2026 年春季的中国模型竞赛表明，
AIME 95\%+ 已经是旗舰模型的“入场门槛”，
推理能力的差异化正在向更难的竞赛级别和真实世界任务转移。

\textbf{长期记忆（LTM）作为 2026 年的“皇冠”能力}——
DeepSeek 的 Engram 论文和 Google 的 Titans 架构
标志着 LLM 正在超越传统的上下文窗口限制，
向真正的“记忆”能力演进。
LTM 不是更长的 context window——
它是一种让模型在运行时自主管理和检索知识的新范式，
可能从根本上改变 AI 系统与人类的交互方式。

\textbf{Agent 生态的成熟}：
从概念验证到生产工具，
代码 Agent 已经改变了软件工程的工作流程。
SWE-bench Verified 接近 80\% 的解决率
标志着 Agent 从「玩具」跨入了「工具」的阶段。
MCP 从 Anthropic 的专有协议成长为
由 Linux 基金会治理、OpenAI/Google/Microsoft 共同采用的行业标准，
Agent 生态系统终于有了自己的“HTTP 时刻”。

\textbf{国产芯片的战略性突破}——
GLM-5 在十万块昇腾 910B 上完成“零 NVIDIA”训练，
标志着中国 AI 芯片独立路线的工程可行性已被验证。
从“被迫替代”到“主动选择”的转变，
可能重塑全球 AI 算力的地缘格局。

\textbf{硬件竞争的全面加剧}——
NVIDIA 的垄断正在被 AMD、Google、AWS、华为等多方力量打破。
Vera Rubin 的 10 倍推理成本降低、
TPU v7 Ironwood 的“推理优先”设计、
Trainium3 的 3nm 制程、
以及 NVLink Fusion 模糊供应商边界的跨界互连，
多供应商策略正从“可选”变为“必选”。

\textbf{端侧 AI 的崛起}——
NPU 硬件的进步使得 3B--7B 模型在手机上成为现实。
Step 3.5 Flash 证明了 MoE 稀疏化可以将数据中心级能力
压缩到消费级 GPU 上运行。
端云混合推理将成为 AI 服务的标准架构。

\textbf{监管框架的形成}——
EU AI Act 的 GPAI 条款已于 2025 年 8 月生效，
MCP 被捐赠给 Linux 基金会标志着 Agent 治理的标准化。
合规成本将成为 AI 商业化的重要变量。

\textbf{开源的“定义性胜利”}——
开源模型不仅在性能上持续逼近闭源模型，
更在 2026 年春季实现了市场份额的质变，
据 OpenRouter 统计，中国开源模型占全球开源 LLM 使用量的近 30\%，
价格战使得“智能太便宜，不值得计量”正在成为现实。
行业可能走向一种混合格局，
基础模型开源，增值服务闭源。

LLM 技术的发展速度令人目眩。
但贯穿始终的核心原则不会改变：
数据的质量决定模型的上限，
计算的规模决定模型的起点，
而人类的智慧——
体现在架构设计、训练策略和工程实现中——
决定了我们能以多高的效率到达那个上限。

\subsection{未解之谜}

在这些令人瞩目的进步之中，
几个\textbf{根本性的开放问题}仍然悬而未决：

\textbf{理解 vs.\ 模拟}。
LLM 是否真正“理解”语言和世界，
还是仅仅是极其精密的统计模式匹配器？
这个问题不仅是哲学思辨——
它有实际的工程后果。
如果 LLM 的能力来自真正的理解，
那么它的泛化能力应该是可靠的；
如果来自统计匹配，
那么在训练分布之外的场景中，
它的行为将不可预测。
2025--2026 年的证据是混合的：
一方面，推理模型在从未见过的数学竞赛题上
展现了令人惊讶的泛化能力；
另一方面，这些模型在简单但不寻常的问题上
仍然会犯令人困惑的错误。

\textbf{对齐的稳固性}。
当前的对齐方法（RLHF、DPO、Constitutional AI）
在多大程度上是“脆弱的表面修饰”
而非“根深蒂固的价值内化”？
如果对齐可以通过巧妙的 jailbreak 提示绕过，
那么它的安全价值就大打折扣。
更深层的问题是：
当模型的能力持续增长，
对齐技术能否保持同步的进步？
Anthropic 的 Responsible Scaling Policy
试图通过能力阈值和安全评估来管理这一风险，
但这依赖于我们能够\textbf{准确评估}模型的危险能力，
而评估能力本身可能落后于模型能力的增长。

\textbf{涌现能力的可预测性}。
LLM 的一些能力（如数学推理、代码生成）
似乎在模型规模跨越某个阈值后“突然”出现~\cite{wei2022emergent}，
这种“涌现”现象既激动人心又令人不安。
如果我们无法预测下一代模型将涌现出什么能力，
那么我们也无法预测它将带来什么风险。
近期的研究（Schaeffer et al., 2024）
对涌现的“突然性”提出了质疑，
认为许多“涌现”可能是度量选择的假象
而非模型能力的真实跃迁，
但这一辩论远未定论。

\subsection{从技术到文明}

这是一个激动人心的时代。
从 Shannon 的信息论到 DeepSeek 的 GRPO，
从 Transformer 的注意力机制到 MoE 的稀疏激活，
从十万块国产芯片上训练出的零 NVIDIA 模型
到 MCP 协议将 Agent 生态推向“HTTP 时刻”：
每一步都是人类对“什么是智能”这个根本问题的持续探索。

2026 年的竞争强度前所未有：
六家中国头部实验室三周内密集发布旗舰模型，
开源与闭源的边界日益模糊，
国产芯片的独立路线从愿景变为现实，
Agent 从“能用”走向“好用”。
AI for Science 正在重塑药物发现和材料科学，
具身智能正在将 AI 从屏幕推向物理世界。
在这场竞赛中，技术进步的速度
已经超越了大多数从业者的预期。

而这段旅程，才刚刚开始。

如果用一句话总结 LLM 训练的核心智慧，那就是：
\textbf{用最简单的目标函数（预测下一个词），
配合足够的数据和计算规模，
让复杂性从简单性中涌现。}
这种“少即是多”的哲学贯穿了整个技术栈：
从 Transformer 的简洁架构到 GRPO 的简单奖励信号，
从 BPE 的朴素合并规则到量化的直观截断思路。
最优雅的工程往往不是最复杂的，而是最简洁的。

让我们回顾整篇文章的核心线索。
LLM 的训练是一个环环相扣的流水线：
\textbf{数据}决定了模型能学到什么（第四章），
\textbf{架构}决定了模型能表达什么（第二、六、七章），
\textbf{训练方法}决定了模型能学多好（第五章），
\textbf{后训练}决定了模型能为人类服务多好（第八章），
\textbf{推理优化}决定了模型能多快多便宜地服务（第九章），
\textbf{推理时计算}决定了模型能“想”多深（第十章）。

每一个环节都蕴含着深刻的思想：
下一个词预测背后是压缩即理解的信息论洞察——
注意力机制背后是让模型自主决定信息流动的设计哲学，
MoE 背后是参数效率与计算效率的解耦思想，
GRPO 背后是让模型通过探索自我发现推理策略的强化学习理念，
量化背后是信息论中“够用的精度”的朴素真理，
Engram 背后是“记忆与计算应该分离”的认知科学洞察。

更宏观地看，2025--2026 年的 LLM 进步
揭示了一个关于\textbf{智能本质}的深刻线索：
智能可能不是一种单一的能力，
而是多种正交维度的组合，
知识存储、逻辑推理、感官感知、
行动规划、社会交互。
LLM 首先在知识存储和语言生成上实现了突破，
然后逐步向推理（o1/R1）、
感知（GPT-4o/Gemini 3 多模态）、
行动（Agent/MCP）和社会交互
（multi-agent 系统）扩展。
这种扩展路径与认知科学对人类智能的理解
有着耐人寻味的呼应。

对于想进入这个领域的读者，
本文的附录提供了完整的训练流水线速查。
对于已经在这个领域工作的实践者，
希望本文的“为什么”视角能帮助你更深入地理解
你每天使用的工具和方法背后的思想。


%% ============================================================
%% NOTE: \appendix command moved to main .tex file (before ch14-appendix)
%% ============================================================
%% BibTeX entries referenced in this chapter
%% (to be added to the main .bib file)
%%
%% @article{deepseekr1,
%%   title   = {{DeepSeek-R1}: Incentivizing Reasoning Capability in {LLMs} via Reinforcement Learning},
%%   author  = {DeepSeek-AI},
%%   journal = {arXiv preprint arXiv:2501.12948},
%%   year    = {2025},
%% }
%%
%% @article{shao2024grpo,
%%   title   = {{DeepSeekMath}: Pushing the Limits of Mathematical Reasoning in Open Language Models},
%%   author  = {Shao, Zhihong and Wang, Peiyi and Zhu, Qihao and others},
%%   journal = {arXiv preprint arXiv:2402.03300},
%%   year    = {2024},
%% }
%%
%% @article{anil2023gemini,
%%   title   = {Gemini: A Family of Highly Capable Multimodal Models},
%%   author  = {Anil, Rohan and Dai, Andrew M and Firat, Orhan and others},
%%   journal = {arXiv preprint arXiv:2312.11805},
%%   year    = {2023},
%% }
%%
%% @article{touvron2023llama,
%%   title   = {{LLaMA}: Open and Efficient Foundation Language Models},
%%   author  = {Touvron, Hugo and Lavril, Thibaut and Izacard, Gautier and others},
%%   journal = {arXiv preprint arXiv:2302.13971},
%%   year    = {2023},
%% }
%%
%% @article{dubey2024llama3,
%%   title   = {The {Llama 3} Herd of Models},
%%   author  = {Dubey, Abhimanyu and Jauhri, Abhinav and Pandey, Abhinav and others},
%%   journal = {arXiv preprint arXiv:2407.21783},
%%   year    = {2024},
%% }
%%
%% @article{hu2021lora,
%%   title   = {{LoRA}: Low-Rank Adaptation of Large Language Models},
%%   author  = {Hu, Edward J and Shen, Yelong and Wallis, Phillip and others},
%%   journal = {arXiv preprint arXiv:2106.09685},
%%   year    = {2021},
%% }
%%
%% @article{dettmers2023qlora,
%%   title   = {{QLoRA}: Efficient Finetuning of Quantized Language Models},
%%   author  = {Dettmers, Tim and Pagnoni, Artidoro and Holtzman, Ari and Zettlemoyer, Luke},
%%   journal = {arXiv preprint arXiv:2305.14314},
%%   year    = {2023},
%% }
%%
%% @article{dao2022flashattention,
%%   title   = {{FlashAttention}: Fast and Memory-Efficient Exact Attention with {IO}-Awareness},
%%   author  = {Dao, Tri and Fu, Daniel Y and Ermon, Stefano and Rudra, Atri and R{\'e}, Christopher},
%%   journal = {Advances in Neural Information Processing Systems},
%%   volume  = {35},
%%   year    = {2022},
%% }
%%
%% @article{gunter2024appleintelligence,
%%   title   = {{Apple Intelligence} Foundation Language Models},
%%   author  = {Gunter, Tom and Wang, Zirui and Wang, Chong and others},
%%   journal = {arXiv preprint arXiv:2407.21075},
%%   year    = {2024},
%% }
%%
%% @article{qwen2025qwen3,
%%   title   = {{Qwen3} Technical Report},
%%   author  = {Qwen Team},
%%   journal = {arXiv preprint},
%%   year    = {2025},
%% }
%%
%% @article{allal2025smollm2,
%%   title   = {{SmolLM2}: When Smol Goes Big --- Data-Centric Training of a Small Language Model},
%%   author  = {{Ben Allal}, Loubna and Lozhkov, Anton and Bakouch, Elie and others},
%%   journal = {arXiv preprint arXiv:2502.02737},
%%   year    = {2025},
%% }
%%
%% --- New references added for 2026-03-12 update ---
%%
%% @article{deepseek2026engram,
%%   title   = {Conditional Memory via Scalable Lookup},
%%   author  = {DeepSeek-AI},
%%   journal = {arXiv preprint},
%%   year    = {2026},
%%   note    = {Published 2026-01-12, Liang Wenfeng co-authored},
%% }
%%
%% @article{smallllm2024toolcalling,
%%   title   = {Small LLMs Outperform Large Models for Agent Tool Calling},
%%   author  = {Various},
%%   journal = {arXiv preprint arXiv:2512.15943},
%%   year    = {2024},
%% }
%%
%% @misc{glm5_2026,
%%   title   = {{GLM-5}: 744B Parameter Model Trained on 100,000 Huawei Ascend 910B},
%%   author  = {Zhipu AI},
%%   year    = {2026},
%%   note    = {Released 2026-02-12, first frontier model with zero NVIDIA dependency},
%% }
%%
%% @misc{kimik25_2026,
%%   title   = {{Kimi K2.5}: 1T MoE Model with Agent Swarm},
%%   author  = {Moonshot AI},
%%   year    = {2026},
%%   note    = {Released 2026-01-27, PARL protocol for multi-agent coordination},
%% }
%%
%% @misc{step35flash_2026,
%%   title   = {{Step 3.5 Flash}: 196B/11B MoE for Consumer Hardware},
%%   author  = {StepFun},
%%   year    = {2026},
%%   note    = {Released 2026-02-02, runs on RTX 4090},
%% }
%%
%% @misc{doubao_seed2_2026,
%%   title   = {{Doubao Seed 2.0 Pro}},
%%   author  = {ByteDance},
%%   year    = {2026},
%%   note    = {Released 2026-02-14, AIME 98.3\%, Codeforces 3020},
%% }
%%
%% @misc{claudecode2_2026,
%%   title   = {{Claude Code 2.0}: Multi-Agent Orchestration for Software Engineering},
%%   author  = {Anthropic},
%%   year    = {2026},
%%   note    = {Released 2026-02-24},
%% }
%%
%% --- New references for ch13 expansion ---
%%
%% @article{peebles2023dit,
%%   title   = {Scalable Diffusion Models with Transformers},
%%   author  = {Peebles, William and Xie, Saining},
%%   journal = {arXiv preprint arXiv:2212.09748},
%%   year    = {2023},
%% }
%%
%% @article{lecun2022path,
%%   title   = {A Path Towards Autonomous Machine Intelligence},
%%   author  = {LeCun, Yann},
%%   journal = {OpenReview preprint},
%%   year    = {2022},
%% }
%%
%% @article{bardes2024vjepa,
%%   title   = {{V-JEPA}: Revisiting Feature Prediction for Learning Visual Representations from Video},
%%   author  = {Bardes, Adrien and Garrido, Quentin and Ponce, Jean and others},
%%   journal = {arXiv preprint arXiv:2404.12451},
%%   year    = {2024},
%% }
%%
%% @article{vjepa2_2025,
%%   title   = {{V-JEPA 2}: Self-Supervised Video Models Enable Understanding, Prediction, and Planning},
%%   author  = {Meta FAIR},
%%   journal = {arXiv preprint},
%%   year    = {2025},
%% }
%%
%% @article{hu2023gaia1,
%%   title   = {{GAIA-1}: A Generative World Model for Autonomous Driving},
%%   author  = {Hu, Anthony and others},
%%   journal = {arXiv preprint arXiv:2309.17080},
%%   year    = {2023},
%% }
%%
%% @article{jumper2021alphafold,
%%   title   = {Highly accurate protein structure prediction with {AlphaFold}},
%%   author  = {Jumper, John and Evans, Richard and Pritzel, Alexander and others},
%%   journal = {Nature},
%%   volume  = {596},
%%   pages   = {583--589},
%%   year    = {2021},
%% }
%%
%% @article{abramson2024alphafold3,
%%   title   = {Accurate structure prediction of biomolecular interactions with {AlphaFold 3}},
%%   author  = {Abramson, Josh and Adler, Jonas and Dunger, Jack and others},
%%   journal = {Nature},
%%   volume  = {630},
%%   pages   = {493--500},
%%   year    = {2024},
%% }
%%
%% @article{watson2023rfdiffusion,
%%   title   = {De novo design of protein structure and function with {RFdiffusion}},
%%   author  = {Watson, Joseph L and Juergens, David and Bennett, Nathaniel R and others},
%%   journal = {Nature},
%%   volume  = {620},
%%   pages   = {1089--1100},
%%   year    = {2023},
%% }
%%
%% @article{dauparas2022proteinmpnn,
%%   title   = {Robust deep learning-based protein sequence design using {ProteinMPNN}},
%%   author  = {Dauparas, Justas and Anishchenko, Ivan and Bennett, Nathaniel and others},
%%   journal = {Science},
%%   volume  = {378},
%%   pages   = {49--56},
%%   year    = {2022},
%% }
%%
%% @article{merchant2023gnome,
%%   title   = {Scaling deep learning for materials discovery},
%%   author  = {Merchant, Amil and Batzner, Simon and Schoenholz, Samuel S and others},
%%   journal = {Nature},
%%   volume  = {624},
%%   pages   = {80--85},
%%   year    = {2023},
%% }
%%
%% @article{price2024gencast,
%%   title   = {GenCast: Diffusion-based ensemble forecasting for medium-range weather},
%%   author  = {Price, Ilan and others},
%%   journal = {Nature},
%%   year    = {2024},
%% }
%%
%% @article{bi2023pangu,
%%   title   = {Accurate medium-range global weather forecasting with 3D neural networks},
%%   author  = {Bi, Kaifeng and Xie, Lingxi and Zhang, Hengheng and others},
%%   journal = {Nature},
%%   volume  = {619},
%%   pages   = {533--538},
%%   year    = {2023},
%% }
%%
%% @article{zhavoronkov2019deep,
%%   title   = {Deep learning enables rapid identification of potent {DDR1} kinase inhibitors},
%%   author  = {Zhavoronkov, Alex and Ivanenkov, Yan A and Aliper, Alex and others},
%%   journal = {Nature Biotechnology},
%%   volume  = {37},
%%   pages   = {1038--1040},
%%   year    = {2019},
%% }
%%
%% @article{brohan2023rt2,
%%   title   = {{RT-2}: Vision-Language-Action Models Transfer Web Knowledge to Robotic Control},
%%   author  = {Brohan, Anthony and Brown, Noah and Carbajal, Justice and others},
%%   journal = {arXiv preprint arXiv:2307.15818},
%%   year    = {2023},
%% }
%%
%% @article{makoviychuk2021isaac,
%%   title   = {{Isaac Gym}: High Performance {GPU}-Based Physics Simulation for Robot Learning},
%%   author  = {Makoviychuk, Viktor and Wawrzyniak, Lukasz and Guo, Yunrong and others},
%%   journal = {NeurIPS Datasets and Benchmarks},
%%   year    = {2021},
%% }
%%
%% @article{kaplan2020scaling,
%%   title   = {Scaling Laws for Neural Language Models},
%%   author  = {Kaplan, Jared and McCandlish, Sam and Henighan, Tom and others},
%%   journal = {arXiv preprint arXiv:2001.08361},
%%   year    = {2020},
%% }
%%
%% @article{hoffmann2022chinchilla,
%%   title   = {Training Compute-Optimal Large Language Models},
%%   author  = {Hoffmann, Jordan and Borgeaud, Sebastian and Mensch, Arthur and others},
%%   journal = {NeurIPS},
%%   year    = {2022},
%% }
%%
%% @article{villalobos2024position,
%%   title   = {Position: Will We Run Out of Data? Limits of {LLM} Scaling Based on Human-Generated Data},
%%   author  = {Villalobos, Pablo and Ho, Anson and Jermyn, Jaime and Lam, Lennart and Weiss, Eric},
%%   journal = {ICML},
%%   year    = {2024},
%% }
%%
%% @article{shumailov2024curse,
%%   title   = {{AI} models collapse when trained on recursively generated data},
%%   author  = {Shumailov, Ilia and Shumaylov, Zakhar and Zhao, Yiren and others},
%%   journal = {Nature},
%%   volume  = {631},
%%   pages   = {755--759},
%%   year    = {2024},
%% }
%%
%% @article{wei2022emergent,
%%   title   = {Emergent Abilities of Large Language Models},
%%   author  = {Wei, Jason and Tay, Yi and Bommasani, Rishi and others},
%%   journal = {Transactions on Machine Learning Research},
%%   year    = {2022},
%% }
%%
%% @article{liu2023llava,
%%   title   = {Visual Instruction Tuning},
%%   author  = {Liu, Haotian and Li, Chunyuan and Wu, Qingyang and Lee, Yong Jae},
%%   journal = {NeurIPS},
%%   year    = {2023},
%% }
%%
%% @article{alayrac2022flamingo,
%%   title   = {Flamingo: a Visual Language Model for Few-Shot Learning},
%%   author  = {Alayrac, Jean-Baptiste and Donahue, Jeff and Luc, Pauline and others},
%%   journal = {NeurIPS},
%%   year    = {2022},
%% }
